% a04_Bh  Appendix: proof of prop:k02:Bh, the one-step bound with the scheme c-number
% local macros (no clash with preamble.tex / preamble_rot.tex / k01 / k02)
\providecommand{\dP}{\,\tfrac{dP}{(2\pi)^3}}
\providecommand{\fl}{\mathrm{fl}}
\providecommand{\Bh}{\ensuremath{(\mathrm{B}^{h})}}
\providecommand{\Kh}{\ensuremath{(\mathrm{K}_{h})}}
\providecommand{\Kpos}{\ensuremath{(\mathrm{K}\text{-}\mathrm{pos})_{N}}}

\section{The one-step bound with the scheme \texorpdfstring{$c$}{c}-number: proof of \texorpdfstring{\cref{prop:k02:Bh}}{Proposition (B\^h)}}\label{app:a04_Bh}

This appendix proves \cref{prop:k02:Bh}. The notation is that of \cref{ssec:k01:setting,ssec:k02:setting}: $k$ is the spacing index, $j$ counts renormalisation steps, $\eta=5^{j-k}$ is the block scale, the diagonal pair at $k$ consists of the $k$-model after $j(k)$ steps and the $(k+1)$-model after $j(k)+1$ steps (both with block scale $\eta_k=5^{j(k)-k}$, the same unit lattice and the same indicator $\chi_{j(k)}$), $S^{(j)}_{N,k}(\cdot;s)=S^{(j)}_k(\cdot;s)-v_{k,j}$ is the normalised action, $\sS_k(h)=\sum_{x\in\Lk\cap\CJ}\eps_k^3h(x)$ and $c_k(h)=z_k\sS_k(h)$. The cube is King's RG cube $\CK$ (\cref{def:s02:Kcube}), $J\ge J_0$ is fixed, and every source is supported at $\ell^\infty$-distance $>1$ from $\R^3\setminus\CJ$.

The proof has three parts. First, on the support of $\chi_j$ the normalised action at the source $(u,h)$ is split exactly into the normalised action at $(u,0)$, the members of the family $\mathcal N_h$ of \cref{hyp:k02:Kh}, and the $h$-linear vacuum constant $\cV^{(j)}_k(h)$ of the three lowest orders (\cref{app:a04:split}). Second, the one-step differences of the first two pieces are bounded by \cref{prop:k01:Kpos} and \cref{hyp:k02:Kh}, uniformly on compact families by a polarisation argument (\cref{app:a04:rows}). Third, $\cV^{(j)}_k(h)+c_k(h)$ is a local functional of $h$: it equals $\bar\beta^{(j)}_k\int h$, with $\bar\beta^{(j)}_k$ the block mean of an infinite-lattice one-point density, up to errors that are super-exponentially small in $1/\eta$, and $\bar\beta^{(j)}_k$ converges along the diagonal pairs at a power rate, the \emph{bulk sequence} (\cref{app:a04:low}). The three pieces are assembled in \cref{app:a04:assembly}.

\subsection{The split of the effective action}\label{app:a04:split}

On the support of $\chi_j$ King's effective action after $j$ steps is \cite[(3.14)]{K1} (cube form \cite[(3.12)]{K2})
\[
S^{(j)}_k(A_j,\phi_j;s)=\tfrac12\ip{A_j}{\Delta^{(j)}A_j}-\ln Z_j+\tfrac12\ip{\phi_j}{\Delta^{(j)}(A^{(j)})\phi_j}-\ln Z_j(A^{(j)})+E_0+P^{(j),1}(A_j,\phi_j;s),
\]
where $P^{(j),1}$ is the sum of the connected renormalised graphs $H$: the Taylor coefficients of King's $E'_j$ \cite[(3.17)]{K1} in the parameters $(e',\lambda',\tau,\theta)$ up to order $\bar n=13$. The sources enter $E'_j$ only through the vertices \cite[(3.20)]{K1} (the curvature source, parameter $\tau$, linear in $u$) and \cite[(3.21)]{K1} (the scalar source, parameter $\theta$, linear in $h$). For a graph $H$ let $n_u(H)$ and $n_h(H)$ be its numbers of curvature and of scalar source vertices, $(\alpha(H),\beta(H))$ its order in $(e^2,\lambda)$, $E(H;A_j,\phi_j;s)$ its value and $E_N(H;\cdot;s)\defeq E(H;\cdot;s)-E(H;0,0;s)$ its field-dependent part. At every field configuration $E(H;\cdot;(u,h))$ is a homogeneous polynomial of degree $n_u(H)$ in $u$ and of degree $n_h(H)$ in $h$, and $n_u(H)+n_h(H)\le\bar n$. The number of graphs is finite and depends on $\bar n$ only.

\begin{definition}[the three families]\label{app:a04:def:families}
Put $P_0\defeq\{(0,0),(0,1),(2,0)\}$ and
\[
\mathcal T_h\defeq\{H:\ n_h(H)=1,\ n_u(H)=0,\ (\alpha(H),\beta(H))\in P_0\},\qquad
\cV^{(j)}_k(h)\defeq\sum_{H\in\mathcal T_h}E(H;0,0;(0,h)),
\]
the \emph{$h$-linear vacuum constant of the three lowest orders}; it is linear in $h$ and does not depend on $u$. The family $\mathcal N_h$ consists of the field-dependent parts $E_N(H;\cdot;(u,h))$ of all graphs with $n_h(H)\ge1$ and of the field-independent values $E(H;0,0;(u,h))$ of all graphs with $n_h(H)\ge1$ and $H\notin\mathcal T_h$. This is the family of \cref{hyp:k02:Kh}; the excepted graphs are exactly those of $\mathcal T_h$, which carry no curvature vertex. In particular the $h$-linear graphs of the orders $P_0$ which do carry curvature vertices (for instance the order-$e^2$ graph in which the seagull vertex joins two curvature legs to the scalar loop through the vertex \cite[(3.21)]{K1}) belong to $\mathcal N_h$.
\end{definition}

\begin{lemma}[the split]\label{app:a04:lem:split}
On the support of $\chi_j$, for every real source $(u,h)$,
\begin{equation}\label{app:a04:eq:split}
S^{(j)}_{N,k}(\cdot;(u,h))\ =\ S^{(j)}_{N,k}(\cdot;(u,0))\ +\ \sum_{T\in\mathcal N_h}T\ +\ \cV^{(j)}_k(h).
\end{equation}
\end{lemma}
\begin{proof}
The quadratic forms, $\ln Z_j$, $\ln Z_j(A^{(j)})$, $E_0$ and the vacuum counterterm $E_1$ inside $V^{(0)}$ are source-free: the sources enter $E'_j$ only through the vertices \cite[(3.20)--(3.21)]{K1}, so they appear in $S^{(j)}_k$ only through $P^{(j),1}$. A graph with $n_h(H)=0$ takes the same value at $(u,h)$ and at $(u,0)$; a graph with $n_h(H)\ge1$ vanishes at $h=0$. Hence
\[
S^{(j)}_k(\cdot;(u,h))-S^{(j)}_k(\cdot;(u,0))=\sum_{n_h(H)\ge1}E(H;\cdot;(u,h))=\sum_{n_h(H)\ge1}E_N(H;\cdot;(u,h))+\sum_{\substack{n_h(H)\ge1\\ H\notin\mathcal T_h}}E(H;0,0;(u,h))+\sum_{H\in\mathcal T_h}E(H;0,0;(u,h)),
\]
and for $H\in\mathcal T_h$ one has $n_u(H)=0$, so $E(H;0,0;(u,h))=E(H;0,0;(0,h))$. Subtracting the source-free constant $v_{k,j}$ from both actions gives \eqref{app:a04:eq:split}.
\end{proof}

\subsection{The field-dependent and the \texorpdfstring{$h$}{h}-nonlinear parts}\label{app:a04:rows}

For a diagonal pair write $D[T]\defeq T|_{(k,j(k))}-T|_{(k+1,j(k)+1)}$ for the one-step difference of a member $T$ of the expansion, a function on the support of $\chi_{j(k)}$ (the two members live on the same unit lattice).

\begin{lemma}[rows at a fixed source]\label{app:a04:lem:rows}
Fix $J$, a family $f_1,\dots,f_n$, $x\in\R^n$ and a real $h$; let $s=(u_x,h)$. There is $k_0$ such that for $k\ge k_0$ (so that $J_*\le j(k)$ and $\eta_k\le\eta_1$), with $j=j(k)$, on the support of $\chi_j$:
\begin{enumerate}[label=\textup{(\alph*)},leftmargin=*]
\item $\big|S^{(j)}_{N,k}(\cdot;(u_x,0))-S^{(j+1)}_{N,k+1}(\cdot;(u_x,0))\big|\le C_3\big(5^{-\gamma j}\eta_k^{-\beta}+\eta_k^{\sigma}\big)$, with the constants of \cref{prop:k01:Kpos} at $s=u_x$;
\item $\big|\sum_{T\in\mathcal N_h}D[T]\big|\le C_h\big(5^{-\gamma'j}\eta_k^{-\beta'}+\eta_k^{\sigma'}\big)|\CJ|$, with $C_h$ the constant of \cref{hyp:k02:Kh} at the source $s$ times the number of members of $\mathcal N_h$, and $\gamma',\beta',\sigma'$ as there.
\end{enumerate}
\end{lemma}
\begin{proof}
(a) is \eqref{eq:k01:Kpos} at $s=u_x$, applied with $j=j(k)$; both members of the pair are admissible since $j(k)+1\le(k+1)-J_*$. (b) is the third clause of \cref{hyp:k02:Kh} for each member of $\mathcal N_h$, which bounds $D[T]$ on the support of $\chi_j$ by $C[5^{-\gamma'j}\eta^{-\beta'}+\eta^{\sigma'}]|\CJ|$, summed over the finitely many members.
\end{proof}

The proposition asks for bounds uniform on a compact family of sources. The families to which it is applied in \cref{sec:k02} are of the form $\Sigma=\{(u_x,h_y):(x,y)\in Y\}$ with $h_y\defeq\sum_jy_jh_j$, $Y\subset\R^{n+p}$ compact, for fixed $f_1,\dots,f_n$ and $h_1,\dots,h_p$; we prove the uniformity for such families.

\begin{lemma}[uniformity on compact families]\label{app:a04:lem:unif}
Let $\Sigma$ be as above and let $W\defeq\{(f_i,0)\}_{i\le n}\cup\{(0,h_j)\}_{j\le p}$. Let $\mathcal S$ be the finite set of sources $\sum_{i=1}^d\epsilon_iw_i$ with $d\le\bar n$, $w_i\in W$ and $\epsilon_i\in\{\pm1\}$. There is $C(Y)<\infty$ such that for every function $T$ of the source which is, at every field configuration, a homogeneous polynomial of degree $d\le\bar n$ in the source $w=(u,h)$,
\[
\sup_{(x,y)\in Y}\big|T(u_x,h_y)\big|\ \le\ C(Y)\,\max_{s'\in\mathcal S}\big|T(s')\big|
\]
pointwise in the fields. This applies to every $E(H)$, $E_N(H)$, to every member of $\mathcal N_h$ and to their one-step differences $D[\cdot]$.
\end{lemma}
\begin{proof}
Let $M$ be the symmetric $d$-linear form with $T(w)=M(w,\dots,w)$. By the polarisation identity, $M(w_1,\dots,w_d)=(2^dd!)^{-1}\sum_{\epsilon\in\{\pm1\}^d}\epsilon_1\cdots\epsilon_d\,T(\sum_i\epsilon_iw_i)$. Expanding $(u_x,h_y)=\sum_ix_i(f_i,0)+\sum_jy_j(0,h_j)$ by multilinearity, $T(u_x,h_y)$ is a sum of at most $(n+p)^d$ monomials in $(x,y)$ of degree $d$, each multiplied by a value $M(w_{a_1},\dots,w_{a_d})$ with $w_{a_i}\in W$, which by polarisation is bounded by $\max_{s'\in\mathcal S}|T(s')|$. The monomials are bounded on the compact set $Y$. Each graph's value is such a polynomial (\cref{app:a04:split}), and so is $D[\cdot]$ of it, since both members of a diagonal pair are evaluated at the same source.
\end{proof}

\begin{corollary}[rows K0 and SRC on a compact family]\label{app:a04:cor:rows}
For $\Sigma$ as above there are $k_0$ and constants $C^{\mathrm{K0}}$, $C^{\mathrm{SRC}}$, depending on the couplings, $J$, $f$, $h_1,\dots,h_p$ and $Y$ only, such that for $k\ge k_0$, $j=j(k)$, every $\sigma=(u_x,h_y)\in\Sigma\cup\{0\}$ and every point of the support of $\chi_j$,
\begin{align}
\big|S^{(j)}_{N,k}(\cdot;(u_x,0))-S^{(j+1)}_{N,k+1}(\cdot;(u_x,0))\big|&\le\theta^{\mathrm{K0}}_k\defeq C^{\mathrm{K0}}\big(5^{-\gamma j}\eta_k^{-\beta}+\eta_k^{\sigma}\big),\label{app:a04:eq:K0}\\
\Big|\sum_{T\in\mathcal N_{h_y}}D[T]\Big|&\le\theta^{\mathrm{SRC}}_k\defeq C^{\mathrm{SRC}}\big(5^{-\gamma'j}\eta_k^{-\beta'}+\eta_k^{\sigma'}\big).\label{app:a04:eq:SRC}
\end{align}
\end{corollary}
\begin{proof}
Write $S^{(j)}_{N,k}(\cdot;(u_x,0))=S^{(j)}_{N,k}(\cdot;0)+\sum_{n_u(H)\ge1,\,n_h(H)=0}E(H;\cdot;(u_x,0))$ (as in the proof of \cref{app:a04:lem:split}, with the roles of $u$ and $h$ exchanged). The first term's one-step difference is \eqref{eq:k01:Kpos} at $s=0$. The second is a polynomial in $x$ without constant term, and its one-step difference is bounded, by \cref{app:a04:lem:unif}, by $C(Y)$ times the maximum over the finitely many curvature sources $s'=(u_{x'},0)\in\mathcal S$ of $|D[S_N(\cdot;(u_{x'},0))]|+|D[S_N(\cdot;0)]|$, each bounded by \cref{app:a04:lem:rows}(a) at $x'$. This gives \eqref{app:a04:eq:K0} with $C^{\mathrm{K0}}$ the largest of the finitely many constants. Likewise $\sum_{T\in\mathcal N_{h_y}}D[T]$ is, at each field configuration, a polynomial in $(x,y)$ (every member has $n_h\ge1$, so it has no $y$-free term), and \cref{app:a04:lem:unif} bounds it by $C(Y)$ times the maximum over $s'\in\mathcal S$ of $|\sum_{T\in\mathcal N_{h'}}D[T](s')|$, each bounded by \cref{app:a04:lem:rows}(b) at the source $s'$; $|\CJ|$ is absorbed into $C^{\mathrm{SRC}}$. The constants $k_0$ are the largest of the finitely many.
\end{proof}

\subsection{The \texorpdfstring{$h$}{h}-linear constant of the three lowest orders}\label{app:a04:low}

\subsubsection{The constant as a fluctuation one-point function}

\begin{lemma}[the vacuum constant is a one-point function]\label{app:a04:lem:V}
Let $\langle\cdot\rangle^{\fl}_{k,j}$ be King's fluctuation measure \cite[(3.17)]{K1} on $\CK$ at zero block fields $A_j=\phi_j=0$: the Gaussian measure of the fluctuation fields with the covariances $G_j(0)$ and $G^{\mathrm{ph}}_{(i)}$ of the $j$ steps, weighted by $e^{V^{(0)}}$ with the vertices \cite[(3.18)--(3.26)]{K1} at zero source, and let $[\,\cdot\,]_{P_0}$ denote the sum of the Taylor coefficients at the orders $P_0$ in $(e^2,\lambda)$. Then
\begin{equation}\label{app:a04:eq:V}
\cV^{(j)}_k(h)=\sum_{x\in\Lk\cap\CJ}\eps_k^3h(x)\,v^{(j)}_k(x),\qquad v^{(j)}_k(x)=-\Big[\big\langle r_x^2-2\sW_{\eps_k}\big\rangle^{\fl}_{k,j}\Big]_{P_0}.
\end{equation}
\end{lemma}
\begin{proof}
At $A_j=\phi_j=0$ the background fields vanish, the quadratic forms and the normalisation constants in \cite[(3.14)]{K1} do not depend on $\theta$, and $P^{(j),1}(0,0;(0,\theta h))$ is the Taylor polynomial of $E'_j=-\ln\int\prod d\mu\,e^{V^{(0)}}$. The $\theta$-derivative of $V^{(0)}$ at $\phi^{(j)}=0$ is the vertex \cite[(3.21)]{K1}, which in unscaled units is $\sum_x\eps_k^3h(x)(r_x^2-2\sW_{\eps_k})$. Differentiating under the finite-dimensional integral, $\partial_\theta E'_j|_{\theta=0}=-\langle\partial_\theta V^{(0)}\rangle^{\fl}_{k,j}$; the graphs of $\mathcal T_h$ are the Taylor coefficients of this derivative at the orders $P_0$ (and at $\tau=0$).
\end{proof}

\subsubsection{The bulk density}

\begin{definition}[King's block spin at zero field on the infinite lattice; the bulk density]\label{app:a04:def:bulk}
On the infinite lattice $\Lk=\eps_k\Z^3$ run King's block-spin averaging at $A=0$ with block factor $5$: the uniform photon average, and for the scalar the weighted average \cite[(3.1)--(3.2)]{K2} with the Dirichlet weight $f_D$ of \cite[(3.8)]{K2} at the finest step and the uniform weight above it (as on the cube, \cref{ssec:k01:setting}). Here $f_D$ is $5\eps_k$-periodic, equal to the constant $(5/|S|)^3$ on the sites of each $5\eps_k$-block whose coordinates relative to the block's corner lie in $S\eps_k$, where $S\subset\{0,\dots,4\}$ is a set of $|S|\in\{3,4\}$ consecutive integers, and zero on the remaining layer(s); its block averages are $1$ and $\norm{f_D}_\infty\le(5/3)^3$. (The printed weight and the corrected weight of \cref{ssec:k01:setting} are both of this form; nothing below depends on which.) Let $Q_j$ be the composite average onto the blocks of side $\eta=5^j\eps_k$, $Q_j^*$ its adjoint for the inner products $\sum\eps_k^3$, $\sum\eta^3$. The fluctuation covariance after $j$ steps is
\[
G_j\defeq\big(-\Delta_{\eps_k}+m^2+\mathfrak a_j\eta^{-2}Q_j^*Q_j\big)^{-1},\qquad \mathfrak a_j=\mathfrak a(1-5^{-2})(1-5^{-2j})^{-1}
\]
\cite[(2.13)]{K1}, \cite[(3.5)]{K2}, and the free propagator splits as $C\defeq(-\Delta_{\eps_k}+m^2)^{-1}=G_j+\Gamma_j$; the photon is the same with $m^2$ replaced by $\lmu^2$, the uniform weight and $A=0$: $D=D_j+\Gamma^{\mathrm{ph}}_j$ with $D_j=\sum_{i<j}G^{\mathrm{ph}}_{(i)}$ \cite[(2.17)]{K1}. Let $v^{(j)}_{k,\infty}(x)$ be the density of \cref{app:a04:lem:V} for this flow on $\Lk$ (the same finitely many graphs, with the free lines $G_j$, $G^{\mathrm{ph}}_{(i)}$ on $\Lk$ and the vertices summed over $\Lk$), and
\[
\beta^{(j)}_k(x)\defeq v^{(j)}_{k,\infty}(x)+z_k,\qquad \bar\beta^{(j)}_k\defeq\eta^{-3}\sum_{x\in\Lk\cap[0,\eta)^3}\eps_k^3\,\beta^{(j)}_k(x).
\]
$\beta^{(j)}_k$ is $\eta\Z^3$-periodic.
\end{definition}

\subsubsection{Decay of the fluctuation covariance at zero field}

\begin{lemma}[decay at $A=0$]\label{app:a04:lem:CT}
Let $f\ge0$ be a weight on $\eps\Z^3$ with $\norm f_\infty<\infty$ and block averages $\eta^{-3}\sum_{x\in B(y)}\eps^3f(x)=1$ over blocks $B(y)$ of side $\eta$; let $(Q_f\phi)(y)\defeq\eta^{-3}\sum_{x\in B(y)}\eps^3f(x)\phi(x)$, so that $Q_f^*Q_f=M_f\mathsf PM_f$ with $\mathsf P$ the uniform block average and $M_f$ multiplication by $f$, and $H\defeq-\Delta_\eps+m^2+b\eta^{-2}Q_f^*Q_f$ with $b\in[b_-,b_+]\subset(0,\infty)$. There are $c_0,\alpha_0>0$, depending on $b_\pm$ and $\norm f_\infty$ only, such that $H\ge c_0\eta^{-2}$ and, for every unit vector $w$ and $|\alpha|\le\alpha_0$, $\norm{e^{\alpha w\cdot x/\eta}H^{-1}e^{-\alpha w\cdot x/\eta}}\le2c_0^{-1}\eta^2$. Consequently the kernel with respect to $\eps^3$ obeys
\[
|H^{-1}(x,y)|\le2c_0^{-1}\eta^2\eps^{-3}e^{-\alpha_0|x-y|_\infty/\eta}.
\]
(King's weights: $f\equiv1$ on the steps above the finest one, $f=f_D$ at the finest step.)
\end{lemma}
\begin{proof}
\emph{Gap.} On a block $B$ write $\bar\phi$ for the uniform mean and $\ell(\phi)\defeq\langle f\phi\rangle_B=(Q_f\phi)(y)$, so $\ell(1)=1$ and $|\ell(\psi)|\le\norm f_\infty\langle|\psi|\rangle_B$. Then $\bar\phi=\ell(\phi)-\ell(\phi-\bar\phi)$, and the discrete Poincar\'e inequality on a cube of side $\eta$ and mesh $\eps$ (uniform in $\eps/\eta$) gives $\langle|\phi|^2\rangle_B\le2(1+\norm f_\infty)^2C_P\eta^2\langle|\nabla\phi|^2\rangle_B+2|\ell(\phi)|^2$. Sum over blocks: $-\Delta_\eps$ dominates the sum of the block Laplacians, and $\sum_y\eta^3|\ell(\phi)|^2=\ip{\phi}{Q_f^*Q_f\phi}$. Hence $\norm\phi^2\le C\max(1,b^{-1})\eta^2\ip{\phi}{H\phi}$. \emph{Combes--Thomas argument} \cite{CT73}. With $\Phi=\alpha w\cdot x/\eta$, $e^{\Phi}(-\Delta_\eps)e^{-\Phi}=-\Delta_\eps+R$ with $\Re\ip{\psi}{R\psi}\ge-\sum_\mu2\eps^{-2}(\cosh(\alpha\eps/\eta)-1)\norm\psi^2\ge-C\alpha^2\eta^{-2}\norm\psi^2$ (the antisymmetric part of the conjugated hopping is skew), and $\norm{e^\Phi M_f\mathsf PM_fe^{-\Phi}-M_f\mathsf PM_f}\le\norm f_\infty^2(e^{\sqrt3\alpha}-1)$ since $\mathsf P$ has range one block. For $\alpha\le\alpha_0$ small, $\Re\,e^\Phi He^{-\Phi}\ge\frac{c_0}2\eta^{-2}$. The kernel bound follows from $|K(x,y)|=\eps^{-6}|\ip{\mathbf 1_x}{K\mathbf 1_y}|$ and $\norm{\mathbf 1_x}^2=\eps^3$.
\end{proof}

\subsubsection{Locality of the vacuum constant}

\begin{lemma}[locality]\label{app:a04:lem:loc}
For every real $h$ supported at $\ell^\infty$-distance $>1$ from $\R^3\setminus\CJ$ and $0\le j\le k-J_0$,
\begin{equation}\label{app:a04:eq:loc}
\Big|\cV^{(j)}_k(h)+z_k\sS_k(h)-\sum_{x\in\Lk}\eps_k^3h(x)\,\beta^{(j)}_k(x)\Big|\ \le\ C\,\norm h_{L^1}\,\eps_k^{-\nu}\,e^{-c/\eta},
\end{equation}
with $C,\nu,c>0$ depending on the couplings and $J$ only. Along the diagonal pairs the right side is $\le C\norm h_{L^1}5^{\nu k}e^{-c5^{\alpha k-1}}$ for both members, which decays faster than any geometric sequence.
\end{lemma}
\begin{proof}
By \cref{app:a04:lem:V}, $v^{(j)}_k(x)$ is a finite sum of Feynman graphs of King's fluctuation measure with one vertex \cite[(3.21)]{K1} at $x$ and lines $G_j(\CK)$, $G^{\mathrm{ph}}_{(i)}(\CK)$ (and their differences and contour sums \cite[(3.27)]{K1} for the averaging vertices), the other vertices summed over $\CK$; the counterterm ${\dm}_{\eps_k}$ and the Wick constant $\sW_{\eps_k}$ are the free-boundary numbers, hence the same on $\CK$ and on $\eps_k\Z^3$. \emph{Image representations.} The region operators are signed sums of free ones over a reflection group: \cite[(3.7)]{K2} for the photon (Neumann) and \cite[(3.9)]{K2} for the scalar (Dirichlet, with the weight $f_D$; ``if we had taken $f=1$, the representation (3.9) would not hold''), with every non-identity image at distance $\ge\dist(x,\partial\CK)+\dist(y,\partial\CK)$ from the pair $(x,y)$; the resulting locality bound is \cite[Prop.~5.1, (5.3)]{K2}. The same holds for every slice and for the lattice differences. \emph{Decay.} Every free fine line obeys \cref{app:a04:lem:CT}: $|G_j(u,w)|\le C\eps_k^{-3}\eta^2e^{-\alpha_0|u-w|/\eta}$, the same for the slices $G_{(i)}$ with $\eta$ replaced by $5^{i+1}\eps_k\le\eta$, and for lattice differences with an extra factor $\eps_k^{-1}$ each. \emph{Graph bound.} Expand every region line as the free line plus its non-identity images, and every vertex sum over $\CK$ as the sum over $\eps_k\Z^3$ minus the sum over the complement. Every term other than the infinite-lattice graph contains an image line $(u,gw)$, with $|u-gw|\ge\dist(u,\partial\CK)+\dist(w,\partial\CK)$, or a vertex outside $\CK$, and the path of fine lines from $x$ to it has length $\ge\dist(x,\partial\CK)-\dist(u,\partial\CK)$, respectively $\ge\dist(x,\partial\CK)$; in both cases the total length along the path is $\ge\dist(x,\partial\CK)\ge1$. The infinite-lattice graphs give $\sum_x\eps_k^3h(x)v^{(j)}_{k,\infty}(x)$, and $z_k\sS_k(h)$ adds $z_k$ per site, which gives $\sum_x\eps_k^3h(x)\beta^{(j)}_k(x)$ (the sums over $\Lk\cap\CJ$ and over $\Lk$ agree, $h$ vanishing outside $\CJ$). In every other term the product of the exponential factors along the path is at most $e^{-\alpha_0/(2\eta)}$ times the remaining half of each factor, which makes the vertex sums converge. The singular and $\eps_k$-dependent prefactors (propagator sup norms, differences, $|{\dm}_{\eps_k}|\le C\eps_k^{-1}$, $|\CJ|$, the finitely many graphs of order $\le e^2$) are at most $C\eps_k^{-\nu}$. Along the diagonal pairs $\eta_k^{-1}=5^{k-j(k)}\ge5^{\alpha k-1}$, so $\eps_k^{-\nu}e^{-c/\eta_k}\le5^{\nu k}e^{-c5^{\alpha k-1}}$; the second member has $\eps_{k+1}=\eps_k/5$ and the same $\eta_k$.
\end{proof}

\subsubsection{Aliasing}

\begin{lemma}[aliasing]\label{app:a04:lem:alias}
Let $\beta$ be an $\eta\Z^3$-periodic function on $\Lk$ with block mean $\bar\beta\defeq\eta^{-3}\sum_{x\in\Lk\cap[0,\eta)^3}\eps_k^3\beta(x)$, and $h\in C^\infty_c(\R^3)$. Then for every integer $r\ge2$
\begin{equation}\label{app:a04:eq:alias}
\Big|\sum_{x\in\Lk}\eps_k^3h(x)\beta(x)-\bar\beta\,\sS_k(h)\Big|\le C_r\norm\beta_\infty\norm{\Delta^rh}_{L^1}\,\eta^{2r},\qquad
\big|\sS_k(h)-\textstyle\int h\big|\le C_r\norm{\Delta^rh}_{L^1}\,\eps_k^{2r}.
\end{equation}
\end{lemma}
\begin{proof}
Let $\Gamma^*\defeq\frac{2\pi}\eta\Z^3$ and $\Lambda^*\defeq\frac{2\pi}{\eps_k}\Z^3\subset\Gamma^*$ be the dual lattices. The periodic function $\beta$ has the finite Fourier expansion $\beta(x)=\sum_{\kappa\in\Gamma^*/\Lambda^*}\hat\beta_\kappa e^{i\kappa\cdot x}$ with $|\hat\beta_\kappa|\le\norm\beta_\infty$ and $\hat\beta_0=\bar\beta$. For $\hat h(p)\defeq\int h(x)e^{-ip\cdot x}dx$ (a Schwartz function) the Poisson summation formula on $\Lk$ gives $\sum_{x\in\Lk}\eps_k^3h(x)e^{i\kappa\cdot x}=\sum_{l\in\Lambda^*}\hat h(l-\kappa)$. For $\kappa=0$ this is $\sS_k(h)=\int h+\sum_{l\in\Lambda^*\setminus0}\hat h(l)$; since $|\hat h(p)|\le\norm{\Delta^rh}_{L^1}|p|^{-2r}$ and $\sum_{l\in\Lambda^*\setminus0}|l|^{-2r}=(\eps_k/2\pi)^{2r}\sum_{n\in\Z^3\setminus0}|n|^{-2r}$, finite for $2r>3$, the second bound follows. For $\kappa\ne0$ the points $l-\kappa$, $l\in\Lambda^*$, $\kappa\in(\Gamma^*/\Lambda^*)\setminus0$, run exactly once through $\Gamma^*\setminus\Lambda^*\subset\Gamma^*\setminus0$, all of modulus $\ge2\pi/\eta$; hence $\sum_{\kappa\ne0}|\hat\beta_\kappa||\sum_l\hat h(l-\kappa)|\le\norm\beta_\infty\norm{\Delta^rh}_{L^1}\sum_{p\in\Gamma^*\setminus0}|p|^{-2r}=C_r\norm\beta_\infty\norm{\Delta^rh}_{L^1}\eta^{2r}$.
\end{proof}

\subsubsection{The free lattice propagator at a rate}

\begin{lemma}[the free lattice propagator approaches the continuum one]\label{app:a04:lem:LP}
For $t>0$ let $C_t(\cdot;m')$ be the kernel of $(-\Delta_t+m'^2)^{-1}$ on $t\Z^3$ with respect to $t^3\sum$, and $C_0(r;m')=(4\pi|r|)^{-1}e^{-m'|r|}$ the continuum kernel. There is an absolute constant $C_*$ such that for all $t>0$, all $m'\ge0$ and all $r\in t\Z^3\setminus0$,
\[
|C_t(r;m')-C_0(r;m')|\le C_*\,t^2|r|^{-3}.
\]
\end{lemma}
\begin{proof}
\emph{Scaling and heat kernels.} $-\Delta_t+m'^2=t^{-2}(-\Delta_1+M^2)$ with $M=m't$, so for $r=tn$, $n\in\Z^3\setminus0$,
\[
C_t(tn;m')=t^{-1}\int_0^\infty e^{-M^2\tau}P_\tau(n)\,d\tau,\qquad C_0(tn;m')=t^{-1}\int_0^\infty e^{-M^2\tau}G_\tau(n)\,d\tau,
\]
where $P_\tau(n)=\prod_\mu q_\tau(n_\mu)$ with $q_\tau(\ell)=e^{-2\tau}I_\ell(2\tau)=(2\pi)^{-1}\int_{-\pi}^{\pi}e^{i\ell\vartheta-2\tau(1-\cos\vartheta)}d\vartheta$ is the heat kernel of $\Delta_1$ on $\Z^3$ and $G_\tau(n)=(4\pi\tau)^{-3/2}e^{-|n|^2/4\tau}$ the continuum one. Since $0<e^{-M^2\tau}\le1$, $|C_t-C_0|\le t^{-1}\int_0^\infty|P_\tau(n)-G_\tau(n)|d\tau$, and it suffices to show $\int_0^\infty|P_\tau(n)-G_\tau(n)|d\tau\le C|n|^{-3}$: the mass enters only as a positive weight $\le1$, so the bound is uniform in $m'\ge0$.

\emph{One-dimensional local limit theorem with a rate.} Let $g_\tau(\ell)=(4\pi\tau)^{-1/2}e^{-\ell^2/4\tau}$. We claim that for $\tau\ge1$ and $\ell\in\Z$, $|q_\tau(\ell)-g_\tau(\ell)|\le C\tau^{-3/2}e^{-c\min(\ell^2/\tau,|\ell|)}$, and that $|q_\tau(\ell)|,|g_\tau(\ell)|\le C\tau^{-1/2}e^{-c\min(\ell^2/\tau,|\ell|)}$ for all $\tau>0$, and moreover $|q_\tau(\ell)|\le C\min(1,\tau^{-1/2})e^{-c\min(\ell^2/\tau,|\ell|)}$. The second bound: for $g$ it is elementary; for $q$ shift the contour of the $\vartheta$-integral to $\Im\vartheta=\xi$ with $\xi\defeq\operatorname{sgn}(\ell)\min(|\ell|/4\tau,1)$ (the integrand is entire and $2\pi$-periodic, so the vertical sides cancel), which produces the factor $e^{-\ell\xi+2\tau(\cosh\xi-1)}\le e^{-|\ell\xi|/2}\le e^{-c\min(\ell^2/\tau,|\ell|)}$, while $\Re[2\tau(1-\cos\vartheta)\cosh\xi]\ge c\tau\vartheta^2$ on $|\vartheta|\le\pi$ gives the factor $C\tau^{-1/2}$; the shifted $\vartheta$-integrand has modulus $\le1$, which gives the factor $1$ in the $\min(1,\tau^{-1/2})$ form. For the first bound, in the region $|\ell|\le\tau^{3/4}$ shift both integrals to $\Im z=\xi\defeq\ell/2\tau$ (so $|\xi|\le\frac12\tau^{-1/4}$); the Gaussian's integral over $|\vartheta|>\pi$ contributes $\le C\tau^{-1}e^{-\ell^2/4\tau-\pi^2\tau}$; on $|\vartheta|\le\pi$ write $\varphi(z)=2(1-\cos z)$, $\psi(z)=z^2$, use $e^{-\tau\varphi}-e^{-\tau\psi}=-\tau(\varphi-\psi)\int_0^1e^{-\tau(s\varphi+(1-s)\psi)}ds$ with $|\varphi-\psi|\le C|z|^4$ on the strip, $\Re\psi=\vartheta^2-\xi^2$ and $\Re\varphi\ge\frac4{\pi^2}\vartheta^2-\xi^2-C\xi^4$; the common prefactor is $e^{-\ell\xi+\tau\xi^2+C\tau\xi^4}=e^{-\ell^2/4\tau}O(1)$ since $\ell^4/\tau^3\le1$ there, and the remaining integral $\tau\int(\vartheta^2+\xi^2)^2e^{-4\tau\vartheta^2/\pi^2}d\vartheta\le C\tau^{-3/2}(1+(\ell^2/\tau)^2)$; altogether $|q-g|\le C\tau^{-3/2}e^{-\ell^2/8\tau}$. In the region $|\ell|>\tau^{3/4}$ one has $\min(\ell^2/\tau,|\ell|)\ge\tau^{1/2}$, and the second bound gives $|q|,|g|\le C\tau^{-1/2}e^{-c\min}\le C\tau^{-3/2}\sup_\tau(\tau e^{-c\tau^{1/2}/2})e^{-c\min/2}$.

\emph{Three dimensions.} Telescope $P_\tau-G_\tau=\sum_\mu\big(\prod_{\nu<\mu}g_\tau(n_\nu)\big)(q_\tau(n_\mu)-g_\tau(n_\mu))\big(\prod_{\nu>\mu}q_\tau(n_\nu)\big)$ and let $\varrho=|n|_\infty\ge1$. For $\tau\ge1$ every factor carries its own decay, so $|P_\tau-G_\tau|\le C\tau^{-5/2}e^{-c\min(\varrho^2/\tau,\varrho)}$; $\int_1^{\varrho}$ gives $\le Ce^{-c\varrho}$ and $\int_\varrho^\infty\tau^{-5/2}e^{-c\varrho^2/\tau}d\tau\le\Gamma(\frac32)(c\varrho^2)^{-3/2}$. For $\tau\le1$, $|P_\tau|\le Ce^{-c\varrho}$ by the second bound in its $\min(1,\tau^{-1/2})$ form (for $\tau\le1$ and $\ell\ne0$, $\min(\ell^2/\tau,|\ell|)=|\ell|$), and $G_\tau(n)\le Ce^{-\varrho^2/8}$. Hence $\int_0^\infty|P_\tau-G_\tau|d\tau\le C|n|^{-3}$.
\end{proof}

At $m'=0$ the lemma is the Green's function asymptotics of simple random walk \cite[Thm.~4.3.1]{LL10}, $C_t(r;0)=t^{-1}g(r/t)$ with $g(n)=(4\pi|n|)^{-1}+O(|n|^{-3})$; the proof above is that argument with the mass carried as a weight. Since $|r|\ge t$, the lemma gives $|C_t(r;m')-C_0(r;m')|\le C_*t^{\gamma}|r|^{-1-\gamma}$ for every $\gamma\in(0,2]$.

\subsubsection{The densities at the three lowest orders}

\begin{lemma}[the bulk density at $A=0$]\label{app:a04:lem:bulkform}
With the notation of \cref{app:a04:def:bulk} on $\eps_k\Z^3$, and $c_\lambda>0$ the combinatorial tadpole constant of King's order-$\lambda$ counterterm (${\dm}^{(\lambda)}=-c_\lambda\lambda\sW_{\eps_k}$, \cite[(3.33)]{K1}), the three orders of $\beta^{(j)}_k=\beta_{(0,0)}+\beta_{(0,1)}+\beta_{(2,0)}$ are
\begin{align*}
\beta_{(0,0)}(x)&=2\,\Gamma_j(x,x),\qquad
\beta_{(0,1)}(x)=-2c_\lambda\lambda\sum_y\eps_k^3G_j(x,y)^2\,\Gamma_j(y,y),\\
\beta_{(2,0)}(x)&=e^2\big[\mathfrak g_x(C,C,C;D)-\mathfrak g_x(G_j,G_j,G_j;D_j)\big]-\mathfrak a_x,
\end{align*}
where $\mathfrak g_x(G_1,G_2,G_3;D')$ is the coefficient of $e^2$ in $\langle r_x^2\rangle$ of the Gaussian matter model with outer lines $G_1,G_3$, internal scalar line $G_2$, photon covariance $D'$, King's bond coupling (exchange and seagull) and the free counterterm ${\dm}^{(e^2)}$; $e^2\mathfrak g_x(C,C,C;D)=z_k$ (the order-$e^2$ coefficient of $\langle|\phi(0)|^2\rangle$ on $\eps_k\Z^3$ with the order-$e^2$ counterterm inserted, \cref{def:s02:sources}); and $\mathfrak a_x$ is the order-$e^2$ contribution of the $A$-dependence of King's averaging, the vertices \cite[(3.24)--(3.26)]{K1} with the contour convention \cite[(2.19), (3.27)]{K1}, to $\langle r_x^2\rangle^{\fl}$.
\end{lemma}
\begin{proof}
Order $1$: $\langle r_x^2\rangle^{\fl}_{(0,0)}=2G_j(x,x)$, so $v_{(0,0)}=-2(G_j(x,x)-\sW_{\eps_k})=2\Gamma_j(x,x)$, and $z_k$ has no order-$1$ part. Order $\lambda$: first-order perturbation theory with the vertices \cite[(3.18)--(3.19)]{K1} at $\phi^{(j)}=0$ gives $\langle r^2_x\rangle^{\fl}_{(0,1)}=-2c_\lambda\lambda\sum_y\eps_k^3G_j(x,y)^2(G_j(y,y)-\sW_{\eps_k})$, and $G_j(y,y)-\sW_{\eps_k}=-\Gamma_j(y,y)$; $z_k$ has no order-$\lambda$ part either. Order $e^2$: the $\theta$-linear vacuum graphs of order $e^2$ built from the original vertices are the exchange graph, the seagull graph and the $\dm$ insertion with the lines $G_j$, $D_j$; the photon tadpole $\langle J\rangle$ vanishes by charge conjugation; the remaining graphs, with the averaging vertices, form $\mathfrak a$. Adding $z_k=e^2\mathfrak g_x(C,C,C;D)$ gives the formula.
\end{proof}

\subsubsection{Order one}

\begin{lemma}[the bulk sequence at order $1$]\label{app:a04:lem:B0}
There is $C$, depending on $\mathfrak a$ and on an upper bound $m_0$ for $m\eta$, such that for all $k$ and $1\le j\le k-1$,
\[
\norm{\beta_{(0,0)}}_\infty\le C\eta^{-1},\qquad\big|\bar\beta^{(j)}_{(0,0),k}-\bar\beta^{(j+1)}_{(0,0),k+1}\big|\le C\eta^{-1}5^{-j/2}.
\]
\end{lemma}

\noindent\emph{Reduction to unit block scale.} Put $t\defeq\eps_k/\eta=5^{-j}$, $m'\defeq m\eta$, $b\defeq\mathfrak a_j$. The map $x\mapsto x/\eta$ carries $-\Delta_{\eps_k}+m^2+\mathfrak a_j\eta^{-2}Q_j^*Q_j$ to $\eta^{-2}(-\Delta_t+m'^2+bQ^*Q)$ on $\Lambda_t\defeq t\Z^3$ with unit blocks; kernels with respect to the lattice volume element scale as $\Gamma_j(x,x')=\eta^{-1}\Gamma_t(x/\eta,x'/\eta)$. Hence $\bar\beta_{(0,0)}=\eta^{-1}F(t;m',b)$ with $F(t;m',b)\defeq2t^3\sum_{x\in\Lambda_t\cap[0,1)^3}\Gamma_t(x,x)$, and the diagonal pair changes $(t,b)\to(t/5,\mathfrak a_{j+1})$ at fixed $m'$, with $|\mathfrak a_j-\mathfrak a_{j+1}|\le C5^{-2j}$. It suffices to prove, uniformly in $m'\in(0,m_0]$ and $b\in[\mathfrak a(1-5^{-2}),\mathfrak a]$:
\begin{itemize}[leftmargin=2em]
\item[(i)] $0\le\Gamma_t(x,x)\le C$;\quad (ii) $|\Gamma_t(x,x')-\Gamma_0(x,x')|\le Ct^{1/2}$ for $x,x'\in\Lambda_t$;\quad (iii) $|\nabla_x\Gamma_0(x,x)|\le C$;\quad (iv) $|\partial_b\Gamma_t(x,x')|\le C$,
\end{itemize}
where $\Gamma_0$ is the continuum block-spin coarse covariance defined below. Then $|F(t)-F(0)|\le2\sup|\Gamma_t-\Gamma_0|+2|t^3\sum\Gamma_0(x,x)-\int_{[0,1)^3}\Gamma_0(u,u)du|\le Ct^{1/2}$ and $|F(t;b)-F(t;b')|\le C|b-b'|$, which gives the lemma.

\begin{proof}[Proof of \cref{app:a04:lem:B0}]
\emph{Fourier representation.} For $t>0$ let $\mathrm{BZ}_t=[-\pi/t,\pi/t)^3$, $\omega_t(p)=\sum_\mu4t^{-2}\sin^2(tp_\mu/2)$, $\hat C_t(p)=(\omega_t(p)+m'^2)^{-1}$; for $t=0$, $\mathrm{BZ}_0=\R^3$, $\omega_0=|p|^2$. Write $p=P+2\pi n$ with $P\in\mathbb T\defeq[-\pi,\pi)^3$ and $n\in N_t$, the set of $t^{-3}$ aliases with $p\in\mathrm{BZ}_t$ ($N_0=\Z^3$). The averaging symbol is $s^f_t(p)\defeq t^3\sum_{x\in\Lambda_t\cap[0,1)^3}f(x)e^{ip\cdot x}$; for the uniform weight $s_t(p)=t^3\sum_{i\in\{0,\dots,t^{-1}-1\}^3}e^{itp\cdot i}$, $s_0(p)=\int_{[0,1)^3}e^{iu\cdot p}du$. With $(Q\phi)(y)=\sum_{x\in B(y)}t^3f(x)\phi(x)$ and $Q^*v(x)=f(x)v(y(x))$ one has $QQ^*=1$ and $(QC_tQ^*)\hat{}\,(P)=W_t(P)\defeq\sum_{n\in N_t}|s^f_t(p)|^2\hat C_t(p)$. Woodbury's identity gives $\Gamma_t=C_tQ^*(b^{-1}+QC_tQ^*)^{-1}QC_t$, hence, by Parseval on $\Z^3$,
\[
\Gamma_t(x,x')=\int_{\mathbb T}\frac{\overline{\hat v^t_x(P)}\,\hat v^t_{x'}(P)}{b^{-1}+W_t(P)}\dP,\qquad \hat v^t_x(P)\defeq\sum_{n\in N_t}s^f_t(p)\hat C_t(p)e^{-ip\cdot x},
\]
and $\Gamma_0$ is \emph{defined} by the same formula at $t=0$ with $s_0$ (the continuum block spin). All sums below converge absolutely.

\emph{Elementary bounds, uniform weight} ($t\in[0,1]$, $P\in\mathbb T$, $n\in N_t$, $p=P+2\pi n$). (E1) $|s_t(p)|=\prod_\mu|t\sin(p_\mu/2)/\sin(tp_\mu/2)|\le\sigma(p)\defeq\prod_\mu\min(1,\pi|\sin(p_\mu/2)|/|p_\mu|)$, using $|\sin(tp_\mu/2)|\ge t|p_\mu|/\pi$ on $\mathrm{BZ}_t$ (and $|s_0|\le\sigma$). If $n_\mu\ne0$ then $|p_\mu|\ge\pi|n_\mu|$, so $\sigma(p)\le\rho(n)\defeq\prod_{\mu:n_\mu\ne0}\frac\pi{2|n_\mu|}$; for $n=0$, $|s_t(P)|\ge(2/\pi)^3$. (E2) $\frac4{\pi^2}|p|^2\le\omega_t(p)\le|p|^2$, so $(|p|^2+m'^2)^{-1}\le\hat C_t(p)\le\frac{\pi^2}4(|p|^2+m'^2)^{-1}$, and $\hat C_t(p)\le\frac34|n|^{-2}$ for $n\ne0$. (E3) For $\kappa<2$, $S_\kappa\defeq\sum_{n\ne0}\rho(n)|n|^{-2+\kappa}<\infty$. (E4) For $p\in\mathrm{BZ}_t$, $s_t(p)=s_0(p)\prod_\mu\mathrm r(itp_\mu)$ with $\mathrm r(z)=z/(e^z-1)$, $|\mathrm r(i\vartheta)|\le\pi/2$ and $|\mathrm r(i\vartheta)-1|\le C|\vartheta|$ for $|\vartheta|\le\pi$, so $|s_t(p)-s_0(p)|\le C\min(1,t|p|)\sigma(p)$; and $|\omega_0-\omega_t|\le t^2|p|^4/12$ gives $|\hat C_t-\hat C_0|\le C\min(t^2,(|p|^2+m'^2)^{-1})\le Ct(|p|^2+m'^2)^{-1/2}$.

\emph{(i).} By (E1)--(E3), $|\hat v^t_x(P)|\le\frac{\pi^2}4(P^2+m'^2)^{-1}+\frac34S_0\le C(P^2+m'^2)^{-1}$, and $W_t(P)\ge|s_t(P)|^2\hat C_t(P)\ge(2/\pi)^6(P^2+m'^2)^{-1}$, so $(b^{-1}+W_t)^{-1}\le(\pi/2)^6(P^2+m'^2)$. The integrand of $\Gamma_t(x,x)$ is $\ge0$ and $\le C(P^2+m'^2)^{-1}$, integrable on $\mathbb T$ uniformly in $m'$.

\emph{(ii).} Put $a_t(n)\defeq s_t(p)\hat C_t(p)$ for $n\in N_t$ and $a_t(n)\defeq0$ otherwise. For $n\in N_t$, by (E4), $|a_t(n)-a_0(n)|\le Ct\sigma(p)(|p|^2+m'^2)^{-1/2}$, which is $\le Ct(P^2+m'^2)^{-1/2}$ for $n=0$ and $\le Ct\rho(n)|n|^{-1}$ for $n\ne0$. For $n\notin N_t$, $|n|\ge(3t)^{-1}$ and $|a_0(n)|\le C\rho(n)|n|^{-2}\le3Ct\rho(n)|n|^{-1}$. The phases are the same for $t$ and $0$. Hence, with (E3) at $\kappa=1$, $|\hat v^t_x-\hat v^0_x|\le Ct(P^2+m'^2)^{-1/2}$ and $|W_t-W_0|\le Ct(P^2+m'^2)^{-1/2}$, and
\[
|\Gamma_t(x,x')-\Gamma_0(x,x')|\le\int_{\mathbb T}\Big[\frac{|\hat v^t_x-\hat v^0_x||\hat v^t_{x'}|+|\hat v^0_x||\hat v^t_{x'}-\hat v^0_{x'}|}{b^{-1}+W_t}+\frac{|\hat v^0_x||\hat v^0_{x'}||W_t-W_0|}{(b^{-1}+W_t)(b^{-1}+W_0)}\Big]\dP\le Ct\int_{\mathbb T}(P^2+m'^2)^{-1/2}\dP\le C't.
\]

\emph{(iii).} $|\nabla_x\hat v^0_x(P)|\le\sum_n|p||a_0(n)|\le C(P^2+m'^2)^{-1/2}+CS_1$, so $|\nabla_x\Gamma_0(x,x)|\le2\int|\nabla\hat v^0_x||\hat v^0_x|(b^{-1}+W_0)^{-1}\le C$. \emph{(iv).} $\partial_b(b^{-1}+W)^{-1}=b^{-2}(b^{-1}+W)^{-2}\le b^{-2}(\pi/2)^{12}(P^2+m'^2)^2$, so $|\partial_b\Gamma_t(x,x')|\le C$.

\emph{The Dirichlet weight.} $f_D$ acts at the finest step only (\cref{app:a04:def:bulk}). At unit block scale, splitting $x=5ti+tm$ with $m\in\{0,\dots,4\}^3$ gives $s^f_t(p)=s_{5t}(p)\,\varphi(tp)$, $\varphi(\vartheta)=\prod_\mu\varphi_1(\vartheta_\mu)$, $\varphi_1(\vartheta)=|S|^{-1}\sum_{m\in S}e^{im\vartheta}$, with $s_{5t}$ the uniform symbol at spacing $5t$; $|\varphi_1|\le1$, $|\varphi_1(\vartheta)-1|\le C|\vartheta|$, and $|\varphi_1(\vartheta)|\ge c_S>0$ for $|\vartheta|\le\pi/5$ (its modulus is $|\sin(|S|\vartheta/2)|/(|S||\sin(\vartheta/2)|)$, whose first zero is at $2\pi/|S|\ge\pi/2$). Repeat the proof with $s^f_t$ in place of $s_t$; only (E1) and (E4) change. Write $p\in\mathrm{BZ}_t$ as $p=q+2\pi\ell/(5t)$ with $q\in\mathrm{BZ}_{5t}$, $\ell\in\{-2,\dots,2\}^3$; $s_{5t}$ is $2\pi/(5t)$-periodic, and $q=P+2\pi n'$ with $n'\in\Z^3$ (as $1/(5t)=5^{j-1}\in\N$). For $\ell=0$: $|s^f_t(p)-s_0(p)|\le|s_{5t}-s_0||\varphi|+|s_0||\varphi(tp)-1|\le C\min(1,t|p|)\sigma(p)$ by (E4) at spacing $5t$, so (ii) goes through verbatim. For $\ell\ne0$: $|s^f_t(p)|\le|s_{5t}(q)|\le\sigma(q)$ by (E1) at spacing $5t$, and $|p_\mu|\ge\pi/(5t)$ for some $\mu$, so $\hat C_t(p)\le25t^2/4$; summing $\sigma(q)$ over $|n'_\mu|\le1/(10t)$ gives $\le C(1+\ln t^{-1})^3$. So these aliases contribute $\le Ct^2(1+\ln t^{-1})^3$ to $\hat v^{f,t}_x$ and to $W^f_t$, while the corresponding $a_0(n)$ lie in the tail $|n|\ge c/t$. The lower bound becomes $W^f_t(P)\ge(2/\pi)^6c_S^6(P^2+m'^2)^{-1}$ (for $j\ge1$, $|tP_\mu|\le\pi/5$). Hence (i)--(iv) hold for $f_D$ with $Ct$ replaced by $Ct(1+\ln t^{-1})^3\le C't^{1/2}$, with the same continuum limit $\Gamma_0$.

\emph{Conclusion.} $\norm{\beta_{(0,0)}}_\infty=2\eta^{-1}\sup_x\Gamma_t(x,x)\le C\eta^{-1}$ by (i), and $|\bar\beta^{(j)}_{(0,0),k}-\bar\beta^{(j+1)}_{(0,0),k+1}|=\eta^{-1}|F(t;m',\mathfrak a_j)-F(t/5;m',\mathfrak a_{j+1})|\le\eta^{-1}\big(|F(t)-F(0)|+|F(0;\mathfrak a_j)-F(0;\mathfrak a_{j+1})|+|F(t/5)-F(0)|\big)\le C\eta^{-1}(t^{1/2}+5^{-2j})$.
\end{proof}

\subsubsection{Order \texorpdfstring{$\lambda$}{lambda}}

\begin{lemma}[the bulk sequence at order $\lambda$]\label{app:a04:lem:B1}
For every $\gamma''\in(0,\frac12)$ there is $C$ such that for all $k$ and $1\le j\le k-1$,
\[
\norm{\beta_{(0,1)}}_\infty\le C\lambda(1+j)^3,\qquad\big|\bar\beta^{(j)}_{(0,1),k}-\bar\beta^{(j+1)}_{(0,1),k+1}\big|\le C\lambda\big(5^{-\gamma''j}+5^{-2j}\big).
\]
\end{lemma}
\begin{proof}
At unit block scale (the reduction of \cref{app:a04:lem:B0}, with $t=5^{-j}$) $\beta_{(0,1)}(x)=-2c_\lambda\lambda\,t^3\sum_yG_t(x,y)^2\Gamma_t(y,y)$, a dimensionless sum. \emph{Decay.} $G_t=C_t-\Gamma_t$; by \cref{app:a04:lem:LP} and $C_t(x,x)\le Ct^{-1}$, $|C_t(x,y)|\le C(|x-y|+t)^{-1}$, and $|\Gamma_t(x,y)|\le\Gamma_t(x,x)^{1/2}\Gamma_t(y,y)^{1/2}\le C$ by \cref{app:a04:lem:B0}(i); so $|G_t(x,y)|\le C(|x-y|+t)^{-1}+C$. \Cref{app:a04:lem:CT} at unit block scale gives $|G_t(x,y)|\le Ct^{-3}e^{-\alpha_0|x-y|_\infty}\le Ct^{-3}e^{-\alpha_0|x-y|/\sqrt3}$; use the first bound for $|x-y|\le R_t\defeq6\sqrt3\,\alpha_0^{-1}\ln t^{-1}$ and the second beyond. With \cref{app:a04:lem:B0}(i) this gives the sup bound: $t^3\sum_{|x-y|\le R_t}\big(|x-y|_t^{-1}+1\big)^2\le C(1+R_t)^3$ and the part $|x-y|>R_t$ is $\le Ct^3$ (as below), so $\norm{\beta_{(0,1)}}_\infty\le C\lambda(1+\ln t^{-1})^3\le C\lambda(1+j)^3$. \emph{Increment.} For $x\ne y$, \cref{app:a04:lem:LP} and \cref{app:a04:lem:B0}(ii) give $|G_t(x,y)-G_0(x,y)|\le Ct^2|x-y|^{-3}+Ct^{1/2}$, where $G_0=C_0-\Gamma_0$. Hence
\[
t^3\sum_y|G_t^2-G_0^2|(x,y)\,\Gamma_t(y,y)\le Ct^3\!\!\sum_{0<|r|\le R_t}\!\!\big(t^2|r|^{-3}+t^{1/2}\big)\big(|r|^{-1}+1\big)+Ct^3G_t(x,x)^2+(\text{the part }|r|>R_t)\le Ct^{1/2}(1+\ln t^{-1})^3,
\]
the part $|r|>R_t$ being $\le Ct^{-6}e^{-2\alpha_0R_t/\sqrt3}\sum_{|r|>R_t}t^3e^{-2\alpha_0(|r|-R_t)/\sqrt3}\le Ct^{-6}t^{12}(1+R_t)^2\le Ct^{3}$; further $|\Gamma_t(y,y)-\Gamma_0(y,y)|\le Ct^{1/2}$; and the Riemann sum of the integrable function $G_0(x,\cdot)^2\Gamma_0(\cdot,\cdot)$ (continuum bounds $|\nabla^iG_0(x,y)|\le C|x-y|^{-1-i}e^{-\delta|x-y|}$, $i\le1$, and \cref{app:a04:lem:B0}(iii)) differs from its integral by $\le Ct^{1/2}$. The cell average in $x$ and the dependence on $b=\mathfrak a_j$ are handled as in \cref{app:a04:lem:B0}; $Ct^{1/2}(1+\ln t^{-1})^3\le C5^{-\gamma''j}$.
\end{proof}

\subsubsection{Order \texorpdfstring{$e^2$}{e\^2}}

\begin{lemma}[the bulk sequence at order $e^2$]\label{app:a04:lem:B2}
For every $\vartheta\in(0,\frac12)$ there are $C$, $C_\vartheta$ and $q\ge0$, depending on $\vartheta$, $\mathfrak a$, $e$, $m$, $\lmu$ and an upper bound $m_0$ for $\max(m,\lmu)\eta$, such that for all $k$ and $1\le j\le k-J_0$,
\[
\norm{\beta_{(2,0)}}_\infty\le Ce^2\,\eta^{-1}(1+j)^{q},\qquad\big|\bar\beta^{(j)}_{(2,0),k}-\bar\beta^{(j+1)}_{(2,0),k+1}\big|\le C_\vartheta e^2\,5^{-\vartheta j}\eta^{-1}.
\]
Both bounds hold for the original-vertex part $e^2[\mathfrak g_x(C,C,C;D)-\mathfrak g_x(G_j,G_j,G_j;D_j)]$ and for the averaging part $\mathfrak a_x$ of \cref{app:a04:lem:bulkform} separately; for the first, the increment bound holds for every $\vartheta\in(0,1)$.
\end{lemma}
% a04_Bh_bulk1  Proof of app:a04:lem:B2, part 1: unit block scale, the order-e^2 density, the sandwich, the self-energy, rows (I)+(II)
\providecommand{\ph}{\mathrm{ph}}
\providecommand{\Mh}{\ensuremath{\hat M}}
\providecommand{\Pih}{\ensuremath{\hat\Pi}}
\providecommand{\cell}{\mathrm{cell}}
\providecommand{\BZ}{\ensuremath{\mathsf B}}
\providecommand{\Bc}{\ensuremath{B^{\mathrm c}}}

The proof of \cref{app:a04:lem:B2} occupies the rest of this subsection. The two parts of $\beta_{(2,0)}$ are treated separately: the original-vertex part $e^2[\mathfrak g_x(C,C,C;D)-\mathfrak g_x(G_j,G_j,G_j;D_j)]$ in \S\S B2.1--B2.6 and the averaging part $\mathfrak a_x$ in \S\S B2.7--B2.10, the last of which also adds the two. Throughout, $L\defeq5$ is the block factor.

\subsubsection*{B2.1 Unit block scale, the order-$e^2$ density, the sandwich}

\emph{Normalisation.} As in the reduction before \cref{app:a04:lem:B0}: $t\defeq\eps_k/\eta=5^{-j}$, $\tau\in\{t,t/5\}$ is the spacing of one member of the diagonal pair at unit block scale, $\Lambda_\tau\defeq\tau\Z^3$, $m'\defeq m\eta$, $\mu'\defeq\lmu\eta$, both in $(0,m_0]$, and $b\defeq\mathfrak a_j\in[\mathfrak a(1-5^{-2}),\mathfrak a]$ with $|\mathfrak a_j-\mathfrak a_{j+1}|\le C5^{-2j}=Ct^2$. Kernels are taken with respect to the counting measure $\tau^3\sum$. $C_\tau\defeq(-\Delta_\tau+m'^2)^{-1}$; $Q$ is the block average over the unit cubes (for the weight $f_D$, $Q_f=QM_f$ with the weight on the $5\tau$-blocks, as in the proof of \cref{app:a04:lem:B0}), $P\defeq Q^*Q$; $\Gamma_\tau\defeq C_\tau Q^*(b^{-1}+QC_\tau Q^*)^{-1}QC_\tau$ is the coarse line and $G_\tau\defeq C_\tau-\Gamma_\tau=(-\Delta_\tau+m'^2+bP)^{-1}$ the fluctuation (outer) line (Woodbury). The photon is componentwise the same with $m'\to\mu'$, the uniform block average and the same $b$ \cite[(2.6), (2.10), (2.16)--(2.17)]{K1}: $D_\tau\defeq(-\Delta_\tau+\mu'^2)^{-1}$, $\Gamma^\ph_\tau$, $G^\ph_\tau\defeq D_\tau-\Gamma^\ph_\tau$; for $f_D$ the photon averaging is unchanged \cite[p.~328]{K2}. The diagonal pair keeps $\eta,m',\mu'$ and the unit cell fixed and changes $(\tau,b,\text{weight})=(t,\mathfrak a_j,f\text{ on }5t\text{-blocks})\to(t/5,\mathfrak a_{j+1},f\text{ on }t\text{-blocks})$. Scaling: $\Gamma_j(x,x')=\eta^{-1}\Gamma_t(x/\eta,x'/\eta)$, and $e^2$ has the dimension of a mass, so the order-$e^2$ density is $e^2$ times a dimensionless unit-scale quantity; all bounds below are uniform in $m',\mu'\in(0,m_0]$, $b$ and $\tau\in(0,1]$ unless said otherwise, and $\eta\le1$ is absorbed into constants.

\begin{lemma}[the order-$e^2$ density]\label{app:a04:lem:dens}
At $\lambda=0$ and zero block fields the matter is Gaussian given $A$ with quadratic form $M_A=-\Delta_A+m'^2+\delta$, $(\Delta_A\phi)(u)=\tau^{-2}\sum_\mu[e^{ie'\tau A_\mu(u)}\phi(u+\tau e_\mu)+e^{-ie'\tau A_\mu(u-\tau e_\mu)}\phi(u-\tau e_\mu)-2\phi(u)]$, where $\delta$ is King's free order-$e^2$ counterterm \cite[(3.33)]{K1}, the same number in the full and in the fluctuation measure \cite[p.~330]{K2}. Writing $M_A=M_0+e'M_1+e'^2(M_2+\delta)+O(e'^3)$,
\[
(M_1\phi)(u)=-i\tau^{-1}\textstyle\sum_\mu[A_\mu(u)\phi(u+\tau e_\mu)-A_\mu(u-\tau e_\mu)\phi(u-\tau e_\mu)],\quad
(M_2\phi)(u)=\tfrac12\sum_\mu[A_\mu(u)^2\phi(u+\tau e_\mu)+A_\mu(u-\tau e_\mu)^2\phi(u-\tau e_\mu)].
\]
With matter covariance $G$ and photon covariance $D'$, the coefficient of $e'^2$ in $\langle r_x^2\rangle=2\E_A[M_A^{-1}(x,x)]$ is $-2(G\,\Pi[G,D']\,G)(x,x)$ with
\[
\Pi[G,D']\defeq\E_{D'}[M_2]+\delta-\E_{D'}[M_1GM_1]\qquad(\text{seagull}+\text{counterterm}-\text{exchange});
\]
explicitly $\E[M_2](u,v)=\frac12\tau^{-3}\sum_\mu[D'(u,u)\mathbf 1_{v=u+\tau e_\mu}+D'(u-\tau e_\mu,u-\tau e_\mu)\mathbf 1_{v=u-\tau e_\mu}]$ and $\E[M_1GM_1](u,v)=-\tau^{-2}\sum_\mu\sum_{\sigma,\sigma'\in\{0,1\}}(-1)^{\sigma+\sigma'}D'(u-\sigma\tau e_\mu,v-(1-\sigma')\tau e_\mu)\,G(u+(1-2\sigma)\tau e_\mu,v-(1-2\sigma')\tau e_\mu)$. Hence, at unit block scale,
\begin{equation}\label{app:a04:eq:dens}
\beta_{(2,0)}(x)+\mathfrak a_x=-2e^2\big[(C_\tau\Pi_CC_\tau)(x,x)-(G_\tau\Pi_GG_\tau)(x,x)\big],\qquad \Pi_C\defeq\Pi[C_\tau,D_\tau],\quad\Pi_G\defeq\Pi[G_\tau,G^\ph_\tau];
\end{equation}
$\Pi_C$ is translation invariant. (Overall signs and the factor $2$ play no role in the bounds and are not tracked.)
\end{lemma}
\begin{proof}
This is the Gaussian algebra of \cref{app:a04:lem:bulkform} and \cref{lem:a05:deriv}: $(M_A)^{-1}=M_0^{-1}-e'M_0^{-1}M_1M_0^{-1}+e'^2[M_0^{-1}M_1M_0^{-1}M_1M_0^{-1}-M_0^{-1}(M_2+\delta)M_0^{-1}]+O(e'^3)$; the determinant does not contribute at order $e'^2$ since $\Tr M_0^{-1}M_1=0$, and $\E_{D'}[M_1]=0$. The kernels follow from $\E[A_\mu(u)A_\nu(v)]=\delta_{\mu\nu}D'(u,v)$. In the full model $G=C_\tau$, $D'=D_\tau$; in King's fluctuation measure at zero block fields $G=G_\tau$, $D'=G^\ph_\tau$ (\cref{app:a04:def:bulk}); the difference of the two is the original-vertex part plus the averaging part, by \cref{app:a04:lem:bulkform}.
\end{proof}

\begin{lemma}[sandwich]\label{app:a04:lem:sandwich}
With $\Gamma\defeq\Gamma_\tau=C_\tau-G_\tau$ (indices dropped),
\begin{equation}\label{app:a04:eq:sandwich}
C\Pi_CC-G\Pi_GG=\underbrace{\Gamma\Pi_CC+C\Pi_C\Gamma}_{(\mathrm I)}\ \underbrace{-\,\Gamma\Pi_C\Gamma}_{(\mathrm{II})}\ +\ \underbrace{G(\Pi_C-\Pi_G)G}_{(\mathrm{III})},\qquad
\Pi_C-\Pi_G=\Sigma^{\mathrm{sg}}[\Gamma^\ph]-\Sigma^{\mathrm{ex}}[\Gamma,D]-\Sigma^{\mathrm{ex}}[G,\Gamma^\ph],
\end{equation}
where $\Sigma^{\mathrm{ex}}[G,D']\defeq\E_{D'}[M_1GM_1]$ and $\Sigma^{\mathrm{sg}}[D']\defeq\E_{D'}[M_2]$. The counterterm $\delta$ cancels in $(\mathrm{III})$; only $(\mathrm I),(\mathrm{II})$ contain the renormalised translation-invariant self-energy $\Pi_C$.
\end{lemma}
\begin{proof}
Expand $G\Pi_CG=(C-\Gamma)\Pi_C(C-\Gamma)$ and add $G(\Pi_C-\Pi_G)G$. The second identity is the bilinearity of $\E_{D'}[M_1GM_1]$ in $(G,D')$ and the linearity of $\E_{D'}[M_2]$ in $D'$, with $\Gamma=C-G$, $\Gamma^\ph=D-G^\ph$.
\end{proof}

\begin{lemma}[derivative placement]\label{app:a04:lem:D}
The current at the bond $(u,\mu)$ is $i[\bar\phi(u)\phi(u+\tau e_\mu)-\bar\phi(u+\tau e_\mu)\phi(u)]=i\tau[\bar\phi\nabla_\mu\phi-(\nabla_\mu\bar\phi)\phi](u)$, $\nabla_\mu$ the forward difference. Hence in every Wick contraction of an exchange term, at each of the two current vertices exactly one of the two matter lines meeting there carries one forward difference, and every term of $(G\Sigma^{\mathrm{ex}}[\cdot,\cdot]G)(x,x)$ is a triangle
\[
\textstyle\sum_{u,v\in\Lambda_\tau}\tau^6\,k_1(x,u)\,k_2(u,v)\,k_3(v,x),\qquad k_1,k_3\in\{G,\nabla G\},\quad k_2=(\text{inner matter line})\times(\text{photon line}),
\]
with arguments shifted by at most one lattice unit. When both differences fall on the inner matter line ($\nabla_u\nabla_v\Gamma(u,v)$ or $\nabla_u\nabla_vG(u,v)$), one summation by parts in $u$ (all sums are finite at fixed $\tau$ and absolutely convergent; $\Lambda_\tau$ has no boundary),
\[
\textstyle\sum_uG(x,u)k(u,v)\nabla_uF(u,v)=-\sum_u\big[(\nabla^*_uG)(x,u)k(u,v)+G(x,u-\tau e_\mu)(\nabla^*_uk)(u,v)\big]F(u,v),
\]
$k$ the photon line, moves one of them. Afterwards every line carries at most one difference. With $w_s(r)\defeq(|r|+\tau)^{-s}$ the exponent bookkeeping is $G\mapsto1$, $\nabla G\mapsto2$, $D\mapsto1$, $\nabla D\mapsto2$, $\Gamma,\nabla\Gamma,\Gamma^\ph,\nabla\Gamma^\ph\mapsto0$ (bounded, \cref{app:a04:lem:Gam} below). Every triangle has slots $(a,b,c)$ with $a,c\le2$, $b\le2$, $a+b+c\le5$; the seagull term $G\Sigma^{\mathrm{sg}}[\Gamma^\ph]G$ is a two-point sum $\sum_u\tau^3G(x,u)\Gamma^\ph(u,u)G(u\pm\tau e_\mu,x)$ with slots $(1,1)$ and no difference.
\end{lemma}
\begin{proof}
The current identity is algebra. The slot list: the inner line $k_2$ is one of $\Gamma\cdot D$, $\Gamma\cdot\nabla D$, $\nabla\Gamma\cdot D$, $G_{\mathrm{in}}\cdot\Gamma^\ph$, $\nabla G_{\mathrm{in}}\cdot\Gamma^\ph$, $G_{\mathrm{in}}\cdot\nabla\Gamma^\ph$ (slot $b\in\{1,2\}$), and $k_1,k_3$ have slots in $\{1,2\}$; e.g.\ $\nabla G\,|\,\Gamma D\,|\,\nabla G$ is $(2,1,2)$, $\nabla G\,|\,\nabla G_{\mathrm{in}}\Gamma^\ph\,|\,G$ is $(2,2,1)$, $G\,|\,\nabla^*D\,\nabla\Gamma\,|\,G$ is $(1,2,1)$. The summation by parts is needed because $\nabla\nabla\Gamma_\tau$ is bounded only by $C(1+\ln\tau^{-1})^3$ (the alias sum $\sum_{N_\tau}\rho(n)$, (E1)--(E3) in the proof of \cref{app:a04:lem:B0}).
\end{proof}

\subsubsection*{B2.2 The renormalised one-loop self-energy}

Let $\Pih_\tau(p)\defeq\sum_r\tau^3\Pi_C(r)e^{-ip\cdot r}$ be the symbol of $\Pi_C$ ($\mathrm{BZ}_\tau$-periodic). By \cref{lem:a05:deriv} (whose $\Pi$ equals $-\tfrac12\Pi_C$ in the normalisation of \cref{app:a04:lem:dens}, since its $K$ is $\tfrac12M_0$ and its $M$ is $\tfrac12(M_A-M_0)$) and the scaling of its part (ii), in the notation of \cref{app:a05:A} (unit lattice, masses $a=\tau m'$, $b=\tau\mu'$),
\begin{equation}\label{app:a04:eq:Pisym}
\Pih_\tau(p)=\tau^{-1}\bar\pi_{\tau m',\tau\mu'}(\tau p),\qquad
\Pih_0(p)\defeq-2\Big[|p|^2\Bc_{m',\mu'}(p)+\tfrac12(\mu'^2-2m'^2)\big(\Bc_{m',\mu'}(p)-\Bc_{m',\mu'}(0)\big)\Big],
\end{equation}
the second being $-2$ times the continuum renormalised one-loop self-energy ($\Bc_{ta,tb}(tq)=t^{-1}\Bc_{a,b}(q)$, \cref{lem:a05:bub}(i), so $\Pih_0(p)=-2\tau^{-1}[|q|^2\Bc_{a,b}(q)+\frac12(b^2-2a^2)(\Bc_{a,b}(q)-\Bc_{a,b}(0))]$ at $q=\tau p$).

\begin{lemma}[self-energy bounds]\label{app:a04:lem:Pi}
Uniformly in $\tau\in(0,1]$ and $m',\mu'\in(0,m_0]$, with $C$ depending on $m_0$ only:
\textup{($\Pi$-0)} $\Pih_\tau(0)=0=\Pih_0(0)$;\quad
\textup{($\Pi$-a)} $|\Pih_\tau(p)|\le C(|p|+m_0)$ for $p\in\mathrm{BZ}_\tau$ and $|\Pih_0(p)|\le C(|p|+m_0)$ on $\R^3$;\quad
\textup{($\Pi$-b)} $|\Pih_\tau(p)-\Pih_0(p)|\le C[\tau^2|p|^3+\tau|p|^2+\tau m_0^2]\le C\tau(1+|p|)(|p|+m_0)$ for $p\in\mathrm{BZ}_\tau$.
\end{lemma}
\begin{proof}
By \cref{lem:a05:dec}, at $q=\tau p\in\BZ$, $-\frac12\bar\pi_{a,b}(q)=\hat q^2B_{a,b}(q)+\frac12(b^2-2a^2)(B_{a,b}(q)-B_{a,b}(0))+\rho_{a,b}(q)$, and by \eqref{app:a04:eq:Pisym} $\Pih_\tau(p)-\Pih_0(p)=-2\tau^{-1}\big\{-\frac12\bar\pi_{a,b}(q)-[|q|^2\Bc+\frac12(b^2-2a^2)(\Bc(q)-\Bc(0))]\big\}$, where
\[
-\tfrac12\bar\pi_{a,b}(q)-\big[|q|^2\Bc+\tfrac12(b^2-2a^2)(\Bc(q)-\Bc(0))\big]=(\hat q^2-|q|^2)\Bc(q)+\hat q^2(B-\Bc)(q)+\tfrac12(b^2-2a^2)\big[(B-\Bc)(q)-(B-\Bc)(0)\big]+\rho(q).
\]
By \cref{lem:a05:el}(c), $|\hat q^2-|q|^2|\le|q|^4/12$, and $\Bc(q)\le1/(8|q|)$ (\cref{lem:a05:bub}(i)), so the first term is $\le|q|^3/96$; $|B-\Bc|\le C_0$ on $\BZ$ (\cref{lem:a05:bub}(ii)) bounds the second by $C_0|q|^2$ and the third by $2C_0(a^2+b^2)$; $|\rho|\le\frac34d_0\hat q^2\le\frac34d_0|q|^2$ (\cref{lem:a05:J}). Multiplying by $2\tau^{-1}$ and substituting $q=\tau p$, $a^2+b^2=\tau^2(m'^2+\mu'^2)\le2\tau^2m_0^2$ gives the first form of ($\Pi$-b); on $\mathrm{BZ}_\tau$, $\tau|p|\le\sqrt3\pi$ gives the second. For ($\Pi$-a): $0<\Bc(q)\le\Bc(0)=(4\pi(a+b))^{-1}$ and $\frac12|b^2-2a^2|\le(a+b)^2$, so $|\Pih_0(p)|\le2\tau^{-1}[|q|/8+(a+b)/(4\pi)]=|p|/4+(m'+\mu')/(2\pi)$; then $|\Pih_\tau|\le|\Pih_0|+C\tau(1+|p|)(|p|+m_0)\le C(|p|+m_0)$ on $\mathrm{BZ}_\tau$. ($\Pi$-0): $\hat q=0$ at $q=0$, and $\rho(0)=0$; $\Pih_\tau(0)=0$ is the statement that King's counterterm \cite[(3.33)]{K1} makes the order-$e^2$ one-particle-irreducible two-point function vanish at zero momentum on the infinite lattice (\cref{lem:a05:deriv}).
\end{proof}

% a04_Bh_bulk1b  Proof of app:a04:lem:B2, part 1b: alias sums with the weight, cell averages, rows (I)+(II)

\subsubsection*{B2.3 The alias sums with the Dirichlet weight}

We use the Fourier representation from the proof of \cref{app:a04:lem:B0}: for $x,x'\in\Lambda_\tau$,
\[
\Gamma_\tau(x,x')=\int_{\mathbb T}\overline{\hat v^\tau_x(P)}\,\hat v^\tau_{x'}(P)\,\Mh_\tau(P)\dP,\qquad \Mh_\tau\defeq(b^{-1}+W_\tau)^{-1},\quad a_\tau(n)\defeq s_\tau(p)\hat C_\tau(p),
\]
with $p=P+2\pi n$, $n\in N_\tau$, and (E1)--(E4) as there; for the weight $f_D$, $s^f_\tau=s_{5\tau}(p)\varphi(\tau p)$, $a^f_\tau$, $\hat v^{f,\tau}_x$, $W^f_\tau$, $\Mh^f_\tau$ are the same expressions with $s^f_\tau$ in place of $s_\tau$ (the Poisson-summation derivation uses only that $f_\tau$ is $1$-periodic, so $Q_fC_\tau Q_f^*$ commutes with $\Z^3$-translations; $Q_fQ_f^*\ne1$ is not used). The continuum object is the same $\Gamma_0$ for both weights. Consequences recorded in that proof: $|\hat v^\tau_x|\le C(P^2+m'^2)^{-1}$, $\Mh_\tau\le\min(b,(\pi/2)^6(P^2+m'^2))$, $|\hat v^\tau_x-\hat v^0_x|,|W_\tau-W_0|\le C\tau(P^2+m'^2)^{-1/2}$, hence $|\Mh_\tau-\Mh_0|\le\Mh_\tau\Mh_0|W_\tau-W_0|\le C\tau(P^2+m'^2)^{3/2}$; for $n\ne0$, $c|n|\le|p|\le C|n|$; all uniform in $m'\in(0,m_0]$.

\begin{lemma}[alias sums]\label{app:a04:lem:Af}
For $\tau\le1$, $P\in\mathbb T$ and $\kappa\in[0,2)$, uniform weight:
\begin{equation}\label{app:a04:eq:star}
\textstyle\sum_{n\in N_\tau}|p|^\kappa|a_\tau(n)|\le C(P^2+m'^2)^{(\kappa-2)/2}+CS_\kappa\le C(P^2+m'^2)^{(\kappa-2)/2}.
\end{equation}
For the weight $f_D$ and $\tau=5^{-j}\le\frac15$, write $p=q+2\pi\ell/(5\tau)$ with $q\in\mathrm{BZ}_{5\tau}$, $\ell\in\{-2,\dots,2\}^3$, $q=P+2\pi n'$, $n'\in N_{5\tau}$ (a bijection $n\leftrightarrow(n',\ell)$); $\ell=0$ are the uniform aliases, $\ell\ne0$ the D-aliases ($|p|_\infty\ge\pi/(5\tau)$, so $\hat C_\tau(p)\le25\tau^2/4$, and $|p|\le\sqrt3\pi/\tau$). Then
\begin{enumerate}[label=\textup{(A\arabic*)},leftmargin=*]
\item $\sum_{\ell\ne0}|p|^\kappa|a^f_\tau(n)|\le C\tau^{2-\kappa}(1+\ln\tau^{-1})^3$;
\item $\sum_{n\in N_\tau}|p|^\kappa|a^f_\tau(n)|\le C(P^2+m'^2)^{(\kappa-2)/2}+CS_\kappa+C\tau^{2-\kappa}(1+\ln\tau^{-1})^3\le C(P^2+m'^2)^{(\kappa-2)/2}$;
\item $W^f_\tau(P)\ge(2/\pi)^6c_S^6(P^2+m'^2)^{-1}$, hence $\Mh^f_\tau\le\min(b,C(P^2+m'^2))$; $|\hat v^{f,\tau}_x|\le C(P^2+m'^2)^{-1}$; $W^f_\tau\le C(P^2+m'^2)^{-1}$;
\item $|\hat v^{f,\tau}_x-\hat v^0_x|+|W^f_\tau-W_0|\le C\tau(P^2+m'^2)^{-1/2}$; $|\Mh^f_\tau-\Mh_0|\le C\tau(P^2+m'^2)^{3/2}$; $|\nabla^{(\tau)}\hat v^{f,\tau}_x-\partial\hat v^0_x|\le C\tau(1+\ln\tau^{-1})^3$, where $\nabla^{(\tau)}$ multiplies the $n$-th term by $d_\tau(p_\mu)\defeq(e^{-i\tau p_\mu}-1)/\tau$ and $\partial$ by $-ip_\mu$.
\end{enumerate}
\end{lemma}
\begin{proof}
\eqref{app:a04:eq:star}: the $n=0$ term is $\le|P|^\kappa\frac{\pi^2}4(P^2+m'^2)^{-1}$ by (E1)--(E2), and the aliases give $\sum_{n\ne0}\rho(n)C|n|^{\kappa}|n|^{-2}=CS_\kappa$ (E3). (A1): each D-alias term is $\le(\sqrt3\pi/\tau)^\kappa(25\tau^2/4)\sigma(q)$ by (E1) at spacing $5\tau$ ($|s^f_\tau(p)|\le|s_{5\tau}(q)|\le\sigma(q)$, as in the Dirichlet-weight paragraph of \cref{app:a04:lem:B0}), and $\sum_{\ell\ne0}\sum_{n'\in N_{5\tau}}\sigma(q)\le124\big[1+\sum_{\emptyset\ne A\subset\{1,2,3\}}\prod_{\mu\in A}\sum_{1\le|n'_\mu|\le1/(10\tau)}\frac\pi{2|n'_\mu|}\big]\le C(1+\ln\tau^{-1})^3$. (A2): the $\ell=0$ terms are bounded as in \eqref{app:a04:eq:star} ($|s^f|\le\sigma(p)\le\rho(n)$); add (A1), and $\tau^{2-\kappa}(1+\ln\tau^{-1})^3\le C_\kappa$. (A3): the $n=0$ term of $W^f$ with $|s^f_\tau(P)|\ge(2/\pi)^3c_S^3$ (valid for $|\tau P_\mu|\le\pi/5$, i.e.\ $\tau\le\frac15$) and (E2); the upper bounds from (A2) at $\kappa=0$. (A4): telescope $a^f_\tau(n)-a_0(n)=(s^f_\tau-s_0)\hat C_\tau+s_0(\hat C_\tau-\hat C_0)$. Uniform aliases: $|s^f_\tau(p)-s_0(p)|\le|s_{5\tau}(p)-s_0(p)||\varphi|+|s_0(p)||\varphi(\tau p)-1|\le C\min(1,\tau|p|)\sigma(p)$ ((E4) at spacing $5\tau$ and $|\varphi_1(\vartheta)-1|\le C|\vartheta|$), and (E4) for $\hat C$, give $|a^f_\tau(n)-a_0(n)|\le C\tau\sigma(p)(|p|^2+m'^2)^{-1/2}$, summed $\le C\tau(P^2+m'^2)^{-1/2}+C\tau S_1$, as for the uniform weight. D-aliases: $\sum|a^f_\tau(n)|\le C\tau^2(1+\ln\tau^{-1})^3$ by (A1) at $\kappa=0$, and their continuum partners $a_0(n)$ with $|n|\ge1/(10\tau)$: $\sum_{|n|\ge N}\rho(n)|n|^{-2}\le CN^{-2}(1+\ln N)^2$. Continuum-only aliases $n\notin N_\tau$ ($|n|\ge(3\tau)^{-1}$): $\le C\tau$ as in \cref{app:a04:lem:B0}(ii). All are $\le C\tau\le C\tau(P^2+m'^2)^{-1/2}$ on $\mathbb T$. For $W^f-W_0$ there is one more factor $|s|\le1$; $|\Mh^f_\tau-\Mh_0|\le\Mh^f_\tau\Mh_0|W^f_\tau-W_0|$. Derivative clause: $|d_\tau(\vartheta)|\le|\vartheta|$, $|d_\tau(\vartheta)+i\vartheta|\le\tau\vartheta^2/2$; so $|\nabla^{(\tau)}\hat v^{f,\tau}_x-\partial\hat v^0_x|\le\sum_{n\in N_\tau}[|a^f_\tau-a_0||p|+|a_0|\tau p^2/2]+\sum_{n\notin N_\tau}|a_0||p|$; the uniform aliases give $C\tau\sum_{n\in N_\tau}\sigma(p)[|p|(|p|^2+m'^2)^{-1/2}+p^2(|p|^2+m'^2)^{-1}]\le C\tau(1+\sum_{N_\tau\setminus0}\rho(n))\le C\tau(1+\ln\tau^{-1})^3$, the D-aliases $C\tau(1+\ln\tau^{-1})^3$ by (A1) at $\kappa=1$, and $\sum_{|n|\ge N}\rho(n)|n|^{-1}\le CN^{-1}(1+\ln N)^2$ the rest.
\end{proof}

\subsubsection*{B2.4 Rows $(\mathrm I)$ and $(\mathrm{II})$: the translation-invariant part with rate}

\begin{lemma}[cell averages]\label{app:a04:lem:cell}
Let $\langle F\rangle_\cell\defeq\tau^3\sum_{x\in\Lambda_\tau\cap[0,1)^3}F(x,x)$. With $Z_\tau(P)\defeq\sum_{n\in N_\tau}|s_\tau|^2\hat C_\tau^3\Pih_\tau(p)$, $W_{2,\tau}\defeq\sum_n|s_\tau|^2\hat C_\tau^2$, $Y_\tau\defeq\sum_n|s_\tau|^2\hat C_\tau^2\Pih_\tau(p)$,
\[
\langle(\mathrm I)\rangle_\cell=2\int_{\mathbb T}\Mh_\tau Z_\tau\dP,\qquad\langle(\mathrm{II})\rangle_\cell=-\int_{\mathbb T}\Mh_\tau^2W_{2,\tau}Y_\tau\dP,
\]
and the same with $s^f_\tau$, $\Mh^f_\tau$ for the weight $f_D$. The continuum ($\tau=0$) expressions are the same with $N_0=\Z^3$, $s_0$, $\hat C_0$, $\Pih_0$.
\end{lemma}
\begin{proof}
$(\Gamma\Pi_CC)(x,x)=\sum_{y,z}\tau^6\Gamma(x,y)\Pi_C(y-z)C(z-x)=\int_{\mathbb T}\Mh_\tau\overline{\hat v_x(P)}\sum_na_\tau(n)\hat C_\tau(p)\Pih_\tau(p)e^{-ip\cdot x}\dP$ (sum the translation-invariant lines in Fourier). Averaging over $x\in\Lambda_\tau\cap[0,1)^3$ with weight $\tau^3$ kills the off-diagonal aliases, $\tau^3\sum_xe^{i(p-p')\cdot x}=\tau^3\sum_xe^{2\pi i(n-n')\cdot x}=\delta_{nn'}$ ($n,n'\in N_\tau$). For $(\mathrm{II})$: $(\Gamma\Pi_C\Gamma)(x,x)=\iint\Mh(P)\Mh(P')\overline{\hat v_x(P)}[\sum_{y,z}\tau^6\hat v_y(P)\Pi_C(y-z)\overline{\hat v_z(P')}]\hat v_x(P')$; the bracket is $\sum_{n,n'}a_\tau(n)\overline{a_\tau(n')}\Pih_\tau(p)\tau^3\sum_ze^{-i(p-p')\cdot z}$, and $\sum_{z\in\Lambda_\tau}\tau^3e^{-i(p-p')\cdot z}$ with $p,p'\in\mathrm{BZ}_\tau$ forces $P=P'$ (as a $\delta$ on $\mathbb T$) and $n=n'$; the cell average of $|\hat v_x|^2$ then gives $W_{2,\tau}$. The weight enters $s$ only, and the Poisson step is weight-free.
\end{proof}

\begin{theorem}[rows $(\mathrm I)+(\mathrm{II})$]\label{app:a04:thm:R}
For both weights, uniformly in $b\in[\mathfrak a(1-5^{-2}),\mathfrak a]$ and $m',\mu'\in(0,m_0]$ (for $f_D$: $\tau=5^{-j}\le\frac15$),
\[
\big|\langle(\mathrm I)+(\mathrm{II})\rangle_{\cell,\tau}-\langle(\mathrm I)+(\mathrm{II})\rangle_{\cell,0}\big|\le C\tau(1+m'^{-1}),\qquad
\big|\langle(\mathrm I)+(\mathrm{II})\rangle_{\cell,0}(b)-\langle(\mathrm I)+(\mathrm{II})\rangle_{\cell,0}(b')\big|\le C|b-b'|.
\]
Hence the diagonal increment obeys $\big|[(\mathrm I)+(\mathrm{II})]_{t/5}(\mathfrak a_{j+1})-[(\mathrm I)+(\mathrm{II})]_t(\mathfrak a_j)\big|\le Ct(1+m'^{-1})+Ct^2\le C't(1+m'^{-1})$.
\end{theorem}
\begin{proof}
\emph{(i) Majorants.} $n=0$ ($p=P\in\mathbb T$): $|s|\le1$, $\hat C\le C(P^2+m'^2)^{-1}$, $|\Pih|\le C(|P|+m_0)$ ($\Pi$-a), $\Mh\le C(P^2+m'^2)$ ((A3) for $f_D$). So $|\Mh Z^{(0)}|\le C(|P|+m_0)(P^2+m'^2)^{-2}$ and $\int_{\mathbb T}\le C\int_0^{\sqrt3\pi}(\varrho+m_0)\varrho^2(\varrho^2+m'^2)^{-2}d\varrho\le C(1+\ln(1/m')+m_0/m')\le C(1+m'^{-1})$. Uniform aliases $n\ne0$: $|s|^2\le\rho(n)^2$, $\hat C\le\frac34|n|^{-2}$, $|\Pih|\le C|n|$, so $|\Mh Z^{(n)}|\le C\rho(n)^2|n|^{-5}$ and $\sum_{n\ne0}\le CS_0$. D-aliases ($f_D$): $|\Mh^fZ^{f,(n)}|\le C\sigma(q)^2(25\tau^2/4)^3C|p|\le C\tau^5\sigma(q)^2$, summed $\le C\tau^5(1+\ln\tau^{-1})^3$ by the count in (A1); likewise $W^f_2$ ($\le C\tau^4(1+\ln\tau^{-1})^3$) and $Y^f$ ($\le C\tau^3(1+\ln\tau^{-1})^3$).
\emph{(ii) Differences at a common $n\in N_\tau$.} Telescope the product $\Mh|s|^2\hat C^3\Pih$ one factor at a time; each factor's difference is $\le C\tau(1+|p|)$ times its majorant: $|s_\tau-s_0|\le C\tau|p|\sigma$ (E4; for $f_D$ and $\ell=0$ the bound of the proof of (A4)), $|\hat C_\tau-\hat C_0|\le C\tau(|p|^2+m'^2)^{-1/2}=C\tau(|p|^2+m'^2)^{1/2}(|p|^2+m'^2)^{-1}$ (E4), ($\Pi$-b), and $|\Mh_\tau-\Mh_0|\le C\tau(P^2+m'^2)^{3/2}=C\tau(P^2+m'^2)^{1/2}(P^2+m'^2)$. For $n=0$ this gives $\int_{\mathbb T}C\tau(1+|P|)(|P|+m_0)(P^2+m'^2)^{-2}\le C\tau(1+m'^{-1})$; for $n\ne0$, $\sum C\tau(1+|n|)\rho(n)^2|n|^{-5}\le C\tau$. D-aliases: bound the $\tau$-term by its majorant of (i) ($O(\tau^3(1+\ln\tau^{-1})^3)$ at worst) and its continuum partner by the tail $|n|\ge1/(10\tau)$: $|\Mh_0Z_0^{(n)}|\le C\rho(n)|n|^{-5}$, $\sum_{|n|\ge1/(10\tau)}\le C\tau^2$.
\emph{(iii) Continuum-only aliases} $n\notin N_\tau$ ($|n|\ge(3\tau)^{-1}$): $|\Mh_0Z_0^{(n)}|\le C\rho(n)|n|^{-5}$, $\sum_{|n|\ge1/(3\tau)}\le C\tau^2$.
\emph{(iv) $(\mathrm{II})$.} The same bookkeeping for $\Mh^2W_2Y$ with the majorants $\Mh^2\le C(P^2+m'^2)^2$, $W_2^{(0)}\le C(P^2+m'^2)^{-2}$, $Y^{(0)}\le C(|P|+m_0)(P^2+m'^2)^{-2}$, alias parts $\le CS_0$ (and the D-alias bounds of (i)); every product of majorants is $\le C[(|P|+m_0)(P^2+m'^2)^{-2}+1]$, integrable with $\int\le C(1+m'^{-1})$, and each factor's difference again costs $C\tau(1+|p|)$.
\emph{(v) $b$-dependence at $\tau=0$.} $\partial_b\Mh=b^{-2}\Mh^2\le C(P^2+m'^2)^2$, so $|\partial_b\langle(\mathrm I)\rangle_0|\le C\int(P^2+m'^2)^2(|P|+m_0)(P^2+m'^2)^{-3}\le C$ and $|\partial_b\langle(\mathrm{II})\rangle_0|\le C\int(P^2+m'^2)^3(|P|+m_0)(P^2+m'^2)^{-4}\le C$ (no infrared loss). The continuum expressions are weight-free (same $\Gamma_0$).
\emph{(vi) Diagonal increment.} Insert $\langle\cdot\rangle_0(\mathfrak a_{j+1})$ and $\langle\cdot\rangle_0(\mathfrak a_j)$; $|\mathfrak a_j-\mathfrak a_{j+1}|\le Ct^2$. Each member is compared with the continuum with its own weight and spacing ($t$ or $t/5$).
\end{proof}

% a04_Bh_bulk2  Proof of app:a04:lem:B2, part 2: the line lemmas for row (III) (coarse lines, free lines with differences, outer lines, power counting)

\subsubsection*{B2.5 The lines of row $(\mathrm{III})$}

\begin{lemma}[coarse lines at spacing $\tau$]\label{app:a04:lem:Gam}
Let $\theta\in(0,1)$. There is $C=C(\theta,\mathfrak a,m_0)$ such that for all $\tau\in(0,1]$ ($\tau\le\frac15$ for $f_D$), $m'\in(0,m_0]$, $b\in[\mathfrak a(1-5^{-2}),\mathfrak a]$, $x,x',h,h'\in\Lambda_\tau$ and every forward difference $\nabla=\nabla^{(\tau)}_\mu$ acting on $x$ or on $x'$:
\begin{enumerate}[label=\textup{(\ensuremath{\Gamma}\arabic*)},leftmargin=*]
\item $0\le\Gamma_\tau(x,x)\le C$, $|\Gamma_\tau(x,x')|\le C$;
\item $|\nabla\Gamma_\tau(x,x')|\le C$;
\item $|\nabla\Gamma_\tau(x+h,x'+h')-\nabla\Gamma_\tau(x,x')|\le C(|h|+|h'|)^\theta$;
\item for $x,x'\in\Lambda_t$ ($t=5^{-j}$, both members of the diagonal pair, each with its own weight): $|\Gamma_{t/5}(x,x')-\Gamma_t(x,x')|\le Ct$ and $|\nabla^{(t/5)}\Gamma_{t/5}(x,x')-\nabla^{(t)}\Gamma_t(x,x')|\le Ct(1+\ln t^{-1})^3$;
\item $|\partial_b\Gamma_\tau(x,x')|\le C$, $|\partial_b\nabla\Gamma_\tau(x,x')|\le C$.
\end{enumerate}
The same holds for the photon coarse line $\Gamma^\ph_\tau$ with $m'\to\mu'$ (uniform weight).
\end{lemma}
\begin{proof}
In the Fourier representation a forward difference in $x$ multiplies the $n$-th term of $\hat v^\tau_x$ by $d_\tau(p_\mu)$, $|d_\tau(\vartheta)|\le|\vartheta|$, and a displacement $x\to x+h$ by $e^{-ip\cdot h}$, $|e^{-ip\cdot h}-1|\le\min(2,|p||h|)\le2^{1-\theta}|p|^\theta|h|^\theta$. Write $(*)_\kappa$ for \eqref{app:a04:eq:star}, resp.\ (A2) for $f_D$, and use $\Mh\le C(P^2+m'^2)$ ((A3) for $f_D$).
($\Gamma$1) is \cref{app:a04:lem:B0}(i), including its Dirichlet-weight paragraph.
($\Gamma$2): $|\nabla\hat v_x|\le(*)_1$, $|\hat v_{x'}|\le(*)_0$: the integrand is $\le C(P^2+m'^2)^{-1/2}$, integrable on $\mathbb T$ uniformly in $m'$ (also at $m'=0$).
($\Gamma$3): displacing $x$: $|\sum_na_\tau d_\tau(p_\mu)(e^{-ip\cdot h}-1)e^{-ip\cdot x}|\le C|h|^\theta(*)_{1+\theta}\le C|h|^\theta(P^2+m'^2)^{(\theta-1)/2}$, times $|\hat v_{x'}|\Mh\le C$; displacing $x'$: $|\hat v_{x'+h'}-\hat v_{x'}|\le C|h'|^\theta(*)_\theta\le C|h'|^\theta(P^2+m'^2)^{(\theta-2)/2}$, times $|\nabla\hat v_x|\Mh\le C(P^2+m'^2)^{1/2}$; in both cases the integrand is $C|h|^\theta(P^2+m'^2)^{(\theta-1)/2}$, integrable since $(\theta-1)/2>-3/2$. For $f_D$ the D-alias part of $(*)_{1+\theta}$ is $C\tau^{1-\theta}(1+\ln\tau^{-1})^3\le C_\theta$ (A2), so the constant stays $\tau$-uniform. ($\theta=1$ fails only because $S_2=\infty$.)
($\Gamma$4), first clause: \cref{app:a04:lem:B0}(ii) at both spacings, each member compared with the same $\Gamma_0$ ($\Lambda_t\subset\Lambda_{t/5}$; for $f_D$ by (A4), each with its own weight): $\le Ct/5+Ct$. Second clause: compare both with $\partial\Gamma_0(x,x')\defeq\int\overline{\sum_na_0(n)(-ip_\mu)e^{-ip\cdot x}}\,\hat v^0_{x'}\Mh_0\dP$. Telescoping $\nabla\hat v^\tau_x\,\hat v^\tau_{x'}\Mh_\tau-\partial\hat v^0_x\,\hat v^0_{x'}\Mh_0$, the first term is $|\nabla\hat v^\tau_x-\partial\hat v^0_x|\,|\hat v^\tau_{x'}|\Mh_\tau\le C\tau(1+\ln\tau^{-1})^3$ (the derivative clause of (A4); for the uniform weight the same computation without D-aliases), and the other two, $|\partial\hat v^0_x||\hat v^\tau_{x'}-\hat v^0_{x'}|\Mh_\tau$ and $|\partial\hat v^0_x||\hat v^0_{x'}||\Mh_\tau-\Mh_0|$, are $\le C\tau(P^2+m'^2)^{-1/2}\cdot C$ by \S B2.3; integrate over $\mathbb T$ and apply at $\tau=t$ and $\tau=t/5$.
($\Gamma$5): $\partial_b\Mh=b^{-2}\Mh^2\le C(P^2+m'^2)^2$, so $|\partial_b\Gamma_\tau|\le C\int(P^2+m'^2)^{-2+2}\dP\le C$ and $|\partial_b\nabla\Gamma_\tau|\le C\int(P^2+m'^2)^{-1/2-1+2}\dP\le C$.
Photon: the same formulas with $m'\to\mu'$, the same $Q$ and $b$ \cite[(2.16)--(2.17)]{K1}.
\end{proof}

\begin{lemma}[the free lattice propagator with differences]\label{app:a04:lem:LPgrad}
For forward lattice differences $\nabla^\alpha_t$ of order $|\alpha|\le2$ there is $C$ such that for all $t>0$, $m'\ge0$ and $r\in t\Z^3\setminus0$,
\[
|\nabla^\alpha_tC_t(r;m')-\partial^\alpha C_0(r;m')|\le C\,t\,|r|^{-2-|\alpha|}.
\]
The same holds for the photon kernel $D_t(\cdot;\mu')$.
\end{lemma}
\begin{proof}
In the notation of the proof of \cref{app:a04:lem:LP}, for $r=tn$: $\nabla^\alpha_tC_t(tn)=t^{-1-|\alpha|}\int_0^\infty e^{-M^2\tau}(\nabla^\alpha_1P_\tau)(n)\,d\tau$ and $\partial^\alpha C_0(tn)=t^{-1-|\alpha|}\int_0^\infty e^{-M^2\tau}(\partial^\alpha G_\tau)(n)\,d\tau$, with $\nabla_1$ the unit-lattice difference, so it suffices to show $\int_0^\infty|\nabla^\alpha_1P_\tau-\partial^\alpha G_\tau|(n)\,d\tau\le C|n|^{-2-|\alpha|}$ ($t^{-1-|\alpha|}|n|^{-2-|\alpha|}=t|r|^{-2-|\alpha|}$). \emph{One dimension.} $\nabla^jq_\tau(\ell)=(2\pi)^{-1}\int_{-\pi}^\pi e^{i\ell\vartheta}(e^{i\vartheta}-1)^je^{-2\tau(1-\cos\vartheta)}d\vartheta$ and $\partial^jg_\tau(\ell)=(2\pi)^{-1}\int_\R e^{i\ell\vartheta}(i\vartheta)^je^{-\tau\vartheta^2}d\vartheta$. Repeat the one-dimensional argument of \cref{app:a04:lem:LP} with these insertions: on the shifted contour $|e^{iz}-1|\le C|z|$ and $|(e^{iz}-1)^j-(iz)^j|\le C|z|^{j+1}$. The term $(iz)^j(e^{-\tau\varphi}-e^{-\tau\psi})$ carries one more factor $|z|^j\le C(\vartheta^2+\xi^2)^{j/2}$, which turns $\tau^{-3/2}$ into $\tau^{-(3+j)/2}$; the term $[(e^{iz}-1)^j-(iz)^j]e^{-\tau\varphi}$ is $O(|z|^{j+1})$ and gives $\tau^{-(2+j)/2}$. Hence, for $\tau\ge1$, $|\nabla^jq_\tau(\ell)-\partial^jg_\tau(\ell)|\le C\tau^{-(2+j)/2}e^{-c\min(\ell^2/\tau,|\ell|)}$, and separately $|\nabla^jq_\tau(\ell)|,|\partial^jg_\tau(\ell)|\le C\tau^{-(1+j)/2}e^{-c\min(\ell^2/\tau,|\ell|)}$ for all $\tau>0$ (the same contour shift). \emph{Three dimensions.} $\nabla^\alpha_1P_\tau(n)=\prod_\mu\nabla^{\alpha_\mu}q_\tau(n_\mu)$; telescope as in \cref{app:a04:lem:LP} with $\varrho=|n|_\infty\ge1$. For $\tau\ge1$ the integrand is $\le C\tau^{-(4+|\alpha|)/2}e^{-c\min(\varrho^2/\tau,\varrho)}$ (one differenced factor at $\tau^{-(2+j)/2}$, the two others at their own decay); $\int_1^\varrho\le Ce^{-c\varrho}$ and $\int_\varrho^\infty\tau^{-(4+|\alpha|)/2}e^{-c\varrho^2/\tau}d\tau\le C\varrho^{-2-|\alpha|}$. For $\tau\le1$ write each forward difference as a sum of at most $2^{|\alpha|}$ shifted values: $|\nabla^\alpha_1P_\tau(n)|\le2^{|\alpha|}\max_{|n'-n|_\infty\le2}|P_\tau(n')|\le Ce^{-c(\varrho-2)}$, since the contour shift at $\xi=\operatorname{sgn}(\ell)$ gives $|q_\tau(\ell)|\le e^{-|\ell|+2\tau(\cosh1-1)}\le Ce^{-|\ell|}$ for $\tau\le1$ (the shifted $\vartheta$-integrand has modulus $\le1$); and $|\partial^\alpha G_\tau(n)|\le C\tau^{-(3+|\alpha|)/2}e^{-\varrho^2/4\tau}$, integrable with $\int_0^1\le Ce^{-\varrho^2/8}$. The finitely many $n$ with $\varrho\le3$ are covered by $\int_0^\infty(1+\tau)^{-(3+|\alpha|)/2}d\tau<\infty$ for the lattice term (bounded for $\tau\le1$ by the same shifted-value argument, $\le C\tau^{-(3+|\alpha|)/2}$ for $\tau\ge1$ by the separate bounds) and by the Gaussian factor for the continuum term. The mass enters only through $0<e^{-M^2\tau}\le1$.
\end{proof}

\begin{lemma}[free and outer lines]\label{app:a04:lem:TIOUT}
Let $w_s(r)\defeq(|r|+\tau)^{-s}$ on $\Lambda_\tau$. Uniformly in $\tau\in(0,1]$ and $m',\mu'\ge0$:
\begin{enumerate}[label=\textup{(T\arabic*)},leftmargin=*]
\item $|\nabla^\alpha C_\tau(r)|\le Cw_{1+|\alpha|}(r)$ for all $r\in\Lambda_\tau$ and $|\alpha|\le2$; the same for $D_\tau$.
\item For $r\in\Lambda_t\setminus0$: $|C_{t/5}(r)-C_t(r)|\le Ct^2|r|^{-3}\le Ctw_2(r)$ and $|\nabla^{(t/5)}C_{t/5}(r)-\nabla^{(t)}C_t(r)|\le Ct|r|^{-3}\le Ctw_3(r)$.
\item For $r,r'\in\Lambda_\tau$ with $|r|\ge8t$, $|r'-r|\le2\sqrt3t$, $\tau\le t$ and $|\alpha|\le1$: $|\nabla^\alpha C_\tau(r')-\nabla^\alpha C_\tau(r)|\le Ctw_{2+|\alpha|}(r)$.
\end{enumerate}
For the outer line $G_\tau=C_\tau-\Gamma_\tau$ (both weights; $b\in[\mathfrak a(1-5^{-2}),\mathfrak a]$) and $i\le1$:
\begin{enumerate}[label=\textup{(O\arabic*)},leftmargin=*]
\item $|\nabla^iG_\tau(x,y)|\le Cw_{1+i}(x-y)$ for $|x-y|\le1$ and $|\nabla^iG_\tau(x,y)|\le C$ for all $x,y$; hence on a ball $|x-y|\le R$, $|\nabla^iG_\tau(x,y)|\le C(R+1)^{1+i}w_{1+i}(x-y)$.
\item With $c\defeq\alpha_0/(2\sqrt3)$ and $R_\tau\defeq(16\sqrt3/\alpha_0)\ln\tau^{-1}+1$: $|\nabla^iG_\tau(x,y)|\le C\tau^4e^{-c|x-y|}$ for $|x-y|\ge R_\tau$, where $\alpha_0=\alpha_0(\mathfrak a,\norm{f}_\infty)>0$ is the constant of \cref{app:a04:lem:CT}, independent of $\tau$, $m'$ and $b$.
\end{enumerate}
The photon lines $D_\tau$, $G^\ph_\tau$ obey the same bounds with $\mu'$.
\end{lemma}
\begin{proof}
(T1): by \cref{app:a04:lem:LPgrad}, $|\nabla^\alpha C_\tau(r)-\partial^\alpha C_0(r)|\le C\tau|r|^{-2-|\alpha|}\le C|r|^{-1-|\alpha|}$ for $r\ne0$ ($\tau\le|r|$), and $|\partial^\alpha C_0(r)|\le C|r|^{-1-|\alpha|}$ for the Yukawa kernel $C_0=e^{-m'|r|}/(4\pi|r|)$, uniformly in $m'\ge0$; at $r=0$, $C_\tau(0)=\int_{\mathrm{BZ}_\tau}\hat C_\tau\dP\le C\tau^{-1}$ and $|\nabla^\alpha C_\tau|\le2^{|\alpha|}\tau^{-|\alpha|}C_\tau(0)$ ($|C_\tau(r)|\le C_\tau(0)$ by positive definiteness). Together $|\nabla^\alpha C_\tau(r)|\le C(|r|+\tau)^{-1-|\alpha|}$. (T2): \cref{app:a04:lem:LP} at both spacings, $|C_{t/5}-C_0|+|C_0-C_t|\le2C_*t^2|r|^{-3}$, and $t^2|r|^{-3}\le4tw_2(r)$ since $|r|\ge t$; \cref{app:a04:lem:LPgrad} likewise for the differences (the coarse forward difference of $C_{t/5}$ at spacing $t/5$ is taken at the points of $\Lambda_t$ with step $t/5$, as the statement says). (T3): a lattice path in $\Lambda_\tau$ from $r$ to $r'$ has $\le3\cdot2\sqrt3t/\tau$ steps; along it $|r''|\ge|r|-2\sqrt3t\ge|r|/2$; each step changes $\nabla^\alpha C_\tau$ by $\le\tau\sup|\nabla^{\alpha+e_\nu}C_\tau(r'')|\le C\tau|r|^{-2-|\alpha|}2^{2+|\alpha|}$ by (T1) at order $|\alpha|+1\le2$; the sum is $\le Ctw_{2+|\alpha|}(r)$. (O1): (T1) and ($\Gamma$1), ($\Gamma$2) of \cref{app:a04:lem:Gam}; on $|r|\le R$, $C\le C(R+1)^sw_s(r)$. (O2): \cref{app:a04:lem:CT} at unit block scale ($\eta=1$, $H=-\Delta_\tau+m'^2+bM_f\mathsf PM_f$, $b\in[\mathfrak a(1-5^{-2}),\mathfrak a]$, block averages of $f$ equal to $1$): $|G_\tau(x,y)|\le2c_0^{-1}\tau^{-3}e^{-\alpha_0|x-y|_\infty}$, so $|\nabla G_\tau|\le C\tau^{-4}e^{-\alpha_0(|x-y|_\infty-\tau)}$; with $|r|_\infty\ge|r|/\sqrt3$, $e^{-\alpha_0|r|/\sqrt3}=e^{-c|r|}e^{-c|r|}\le e^{-cR_\tau}e^{-c|r|}\le\tau^8e^{-c|r|}$ for $|r|\ge R_\tau$. Photon: (T1)--(T3) with $\mu'$ (\cref{app:a04:lem:LP,app:a04:lem:LPgrad} hold for every mass), (O1)--(O2) with the uniform weight.
\end{proof}

\begin{lemma}[power counting on confined triangles]\label{app:a04:lem:PC}
Let $\tau\in(0,1]$, $R\ge2$, $x\in\Lambda_\tau$, $a,b,c\in[0,3]$ with $\Sigma\defeq a+b+c\le6$, and $E_R\defeq\{(u,v)\in\Lambda_\tau^2:|x-u|\le R,|x-v|\le R\}$. Then, with $C$ depending only on an upper bound for $\Sigma$,
\begin{equation}\label{app:a04:eq:PC}
\tau^6\sum_{(u,v)\in E_R}w_a(x-u)w_b(u-v)w_c(v-x)\le C(1+\ln(R/\tau))^4R^{6-\Sigma}.
\end{equation}
Two-point version: $\tau^3\sum_{|x-u|\le R}w_a(x-u)w_c(u-x)\le C(1+\ln(R/\tau))R^{(3-a-c)_+}\tau^{-(a+c-3)_+}$. Tail version: for $s\le3$ and $\lambda>0$, $\tau^3\sum_ue^{-\lambda|x-u|}w_s(x-u)\le C_\lambda(1+\ln\tau^{-1})$.
\end{lemma}
\begin{proof}
\emph{Coincident terms.} $u=x$: $\tau^3w_a(0)=\tau^{3-a}$, and the $v$-sum is $\tau^3\sum_{|v-x|\le R}w_{b+c}(v-x)\le C(1+\ln(R/\tau))R^{(3-b-c)_+}\tau^{-(b+c-3)_+}$ (shells); the product is $\le C(1+\ln(R/\tau))R^{6-\Sigma}$ (if $b+c\ge3$ the $\tau$-exponent is $6-\Sigma\ge0$; if $b+c<3$, $\tau^{3-a}R^{3-b-c}\le R^{6-\Sigma}$). $v=x$ and $u=v$: the same. \emph{Non-coincident terms.} For $r\in\Lambda_\tau\setminus0$ and $r'$ in the sub-cube $Q'_r$ of side $\tau/4$ centred at $r$, $|r'-r|\le\sqrt3\tau/8$, so $|r'|\le2(|r|+\tau)$ and $\tau^3w_s(r)\le64\cdot2^s\int_{Q'_r}|r'|^{-s}dr'$. Applied to the three factors with $u'\in Q'_u$, $v'\in Q'_v$ (so $|(u'-v')-(u-v)|\le\sqrt3\tau/4$ and $|u'-v'|\le2(|u-v|+\tau)$): for distinct lattice points the sub-cubes are disjoint and every pairwise distance among $x,u',v'$ is $\ge\tau-\sqrt3\tau/4\ge\tau/2$, whence
\[
\tau^6\sum_{E_R,\ \text{non-coinc.}}\prod w\le64^2\,2^\Sigma\iint_\Omega|x-u'|^{-a}|u'-v'|^{-b}|v'-x|^{-c}\,du'dv',\quad\Omega\defeq\{|x-u'|,|x-v'|\le R+1,\ \text{all pairwise distances}\ge\tau/2\}.
\]
Dyadic shells in $\varrho\defeq\max(|x-u'|,|x-v'|)\in[\tau/2,R+1]$: scaling $(u'-x,v'-x)=2^i(u'',v'')$ on the shell $2^i\le\varrho<2^{i+1}$ gives $2^{i(6-\Sigma)}J_i$, $J_i\defeq\int_{1\le\varrho''<2,\ \text{distances}\ge\delta_i\defeq\tau/2^{i+1}}|u''|^{-a}|u''-v''|^{-b}|v''|^{-c}$. On the shell one of $|u''|,|v''|$ is $\ge1$, say $|u''|\ge1$, so $|u''|^{-a}\le1$; the remaining $\int_{|v''|\le2}|v''|^{-c}|u''-v''|^{-b}dv''$ with $b,c\le3$ and the two singular points at distance $\ge1$ is $\le C(1+\ln(2^i/\tau))^2$ (a slot equal to $3$ costs one logarithm, a slot $<3$ none), and the $u''$-integral over $1\le|u''|\le2$ is bounded. So $J_i\le C(1+\ln(R/\tau))^2$ and $\sum_{i:\tau/4\le2^i\le R+1}2^{i(6-\Sigma)}\le C(1+\ln(R/\tau))R^{6-\Sigma}$ ($\Sigma\le6$). Total $\le C(1+\ln(R/\tau))^4R^{6-\Sigma}$. Two-point and tail versions: shells in one variable.
\end{proof}

\begin{lemma}[one inner sum]\label{app:a04:lem:Tp}
For $b,c\in[0,3]$ and $u,x\in\Lambda_\tau$: $\tau^3\sum_{v\in\Lambda_\tau,|v-x|\le R}w_b(u-v)w_c(v-x)\le C(1+\ln(R/\tau))\cdot\{w_{b+c-3}(u-x)$ if $b+c>3$; $1$ if $b+c=3$; $R^{3-b-c}$ if $b+c<3\}$.
\end{lemma}
\begin{proof}
Split into $|v-x|\le|u-x|/2$ (there $w_b(u-v)\le2^bw_b(u-x)$), $|v-u|\le|u-x|/2$ (there $w_c(v-x)\le2^cw_c(u-x)$) and the rest (there $|v-u|\ge|u-x|/2$ and $|v-x|\ge|u-x|/2$, so $w_bw_c\le Cw_{b+c}(v-x)$ up to the roles), and sum in shells.
\end{proof}

% a04_Bh_bulk3  Proof of app:a04:lem:B2, part 3: row (III) with rate, the sup bound, the original-vertex part assembled

\subsubsection*{B2.6 Row $(\mathrm{III})$ with rate, the sup bound, and the original-vertex part}

\begin{theorem}[row $(\mathrm{III})$]\label{app:a04:thm:III}
Let $\cT_\tau(b)\defeq\tau^3\sum_{x\in\Lambda_\tau\cap[0,1)^3}\big(G(\Pi_C-\Pi_G)G\big)_\tau(x,x)$ (all lines at spacing $\tau$ with block coupling $b$ and the weight of that member). For every $\theta<1$ there is $C_\theta$, depending on $\theta$, $\mathfrak a$, $m_0$ and $\norm{f}_\infty$ only, such that for all $t=5^{-j}\le1$ ($\le\frac15$ for $f_D$) and $m',\mu'\in(0,m_0]$:
\[
\textup{(III-sup)}\ \ \sup_x\big|\big(G(\Pi_C-\Pi_G)G\big)_t(x,x)\big|\le C(1+\ln t^{-1})^{14};\qquad
\textup{(III-inc)}\ \ \big|\cT_{t/5}(\mathfrak a_{j+1})-\cT_t(\mathfrak a_j)\big|\le C_\theta t^\theta.
\]
\end{theorem}
\begin{proof}
By \cref{app:a04:lem:sandwich,app:a04:lem:D}, $(G(\Pi_C-\Pi_G)G)(x,x)$ is a finite sum (independent of $\tau$) of triangle terms $T_\tau(x)=\tau^6\sum_{u,v}k_1(x,u)k_2(u,v)k_3(v,x)$ and of two-point (seagull) terms, with $k_1,k_3\in\{G_\tau,\nabla G_\tau\}$ (slots $a,c\in\{1,2\}$), $k_2\in\{\Gamma D,\Gamma\nabla D,\nabla\Gamma D,G_{\mathrm{in}}\Gamma^\ph,\nabla G_{\mathrm{in}}\Gamma^\ph,G_{\mathrm{in}}\nabla\Gamma^\ph\}$ (slot $b\in\{1,2\}$), $a+b+c\le5$, arguments shifted by at most one lattice unit. Put $R\defeq R_{t/5}+5\le C(1+\ln t^{-1})$ and split each outer line $k_i=k_i^{\mathrm{in}}+k_i^{\mathrm{out}}$, $k^{\mathrm{in}}\defeq k\mathbf 1_{|r|\le R}$. \emph{Sup bounds} (L1): $|k_1|\le C(R+1)^aw_a$, $|k_3|\le C(R+1)^cw_c$ on $|r|\le R$ and $\le C\tau^4e^{-c|r|}$ on $|r|\ge R-5$ (O1), (O2); $|k_2|\le Cw_b$ (($\Gamma$1), ($\Gamma$2), (T1), (O1)), with $C$ replaced by $C(2R+1)^b$ for an inner $G_{\mathrm{in}}$, $\nabla G_{\mathrm{in}}$ on the confined region $|u-v|\le2R$ (a polylogarithm, absorbed below); where an outer line is ``out'', such an inner line is bounded by $Cw_b(u-v)$ on $|u-v|\le1$ and by $C$ beyond (slot $0$). The weight enters only through \cref{app:a04:lem:Gam,app:a04:lem:TIOUT}; the kernels are evaluated, never their weights.

\emph{Step 0 (sup bound).} Core (both outer lines ``in''): \cref{app:a04:lem:PC} with $\Sigma\le5$ gives $\le C(R+1)^{a+c}(2R+1)^b(1+\ln(R/\tau))^4R^{6-\Sigma}\le C(1+\ln t^{-1})^{13}$. Terms with $k_1^{\mathrm{out}}$: $\le C\tau^4\sum_u\tau^3e^{-c|x-u|}[\tau^3\sum_vw_b(u-v)|k_3(v,x)|]$; by \cref{app:a04:lem:Tp} (the $v$-sum over the ball $|v-x|\le R$) and the $k_3^{\mathrm{out}}$ tail, the bracket is $\le C(R+1)^{c+2}(1+\ln\tau^{-1})(1+w_{(b+c-3)_+}(u-x))+C\tau^4(1+\ln\tau^{-1})$ (the case $b+c<3$ of \cref{app:a04:lem:Tp}, with slot $0$ for an inner $G_{\mathrm{in}}$ beyond $|u-v|\le1$, gives $R^{3-b-c}\le R^2$), and the tail version of \cref{app:a04:lem:PC} gives $\le C\tau^4(1+\ln\tau^{-1})^{q}$. Same for $k_3^{\mathrm{out}}$. Seagull terms: the two-point version of \cref{app:a04:lem:PC}. Hence (III-sup) with exponent $14$, uniformly in $m',\mu'$: no line needs mass decay, the outer lines decaying at the $m'$-independent rate $c$.

\emph{Step 1 (decomposition of the increment).} With the cells $c(x)\defeq x+(t/5)\{0,\dots,4\}^3$ (they tile $\Lambda_{t/5}$, and the fine unit cell is the disjoint union of the $c(x)$, $x\in\Lambda_t\cap[0,1)^3$), for each triangle term
\[
\cT_{t/5}(\mathfrak a_j)-\cT_t(\mathfrak a_j)=\cD+\cC,\quad
\cD\defeq t^9\!\!\sum_{x,u,v\in\Lambda_t}\!\Big[5^{-9}\!\!\sum_{x'\in c(x),u'\in c(u),v'\in c(v)}\!\!F_{t/5}(x',u',v')-F_{t/5}(x,u,v)\Big],\quad
\cC\defeq t^9\!\!\sum_{x,u,v\in\Lambda_t}\![F_{t/5}-F_t](x,u,v),
\]
$F_\tau(x,u,v)\defeq k_1k_2k_3$ at spacing $\tau$ ($x$ restricted to the unit cell). Near-diagonal set $\cN\defeq\{\min(|x-u|,|u-v|,|v-x|)<8t\}$. (For $f_D$ the fine member's weight lives on the $t$-blocks, which are exactly the cells $c(x)$, the coarse member's on the $5t$-blocks.)

\emph{Step 2 ($\cN$).} Each of the four sums restricted to $\cN$ (fine triples whose coarse parent is in $\cN$ lie in $\cN'\defeq\{\min<8t+2\sqrt3t\}$) is bounded by (L1) alone. Case $|x-u|<8t$: $\le17^3$ values of $u$ (resp.\ $(17\cdot5)^3$ fine ones), $|F|\le C(R+1)^{a+c}\tau^{-a}w_b(u-v)w_c(v-x)$; \cref{app:a04:lem:Tp} ($R$-confinement by $k_3$, or its tail) gives $\tau^3\sum_v\le C(1+\ln\tau^{-1})R^3t^{-(b+c-3)_+}$, so the contribution is $\le Ct^3t^{-a}t^{-(b+c-3)_+}(1+\ln t^{-1})^{8}\le Ct(1+\ln t^{-1})^{8}$, since $a\le2$ and $a+b+c\le5$ give $3-a-(b+c-3)_+\ge\min(3-a,6-\Sigma)\ge1$. Case $|v-x|<8t$: symmetric. Case $|u-v|<8t$: $\le17^3$ values of $v$ per $u$, $|F|\le C(R+1)^{a+c}w_a(x-u)t^{-b}w_c(v-x)$ with $w_c(v-x)\le2^cw_c(u-x)$ if $|u-x|\ge16t$ and $\le t^{-c}$ otherwise: $\le Ct^{3-b}(R+1)^{a+c}[t^3\sum_{|x-u|\le R}w_{a+c}(x-u)+t^3\cdot33^3t^{-a-c}]\le Ct(1+\ln t^{-1})^{8}$ ($b\le2$, $a+c\le4$, $\Sigma\le5$). Two-point terms: the same with one fewer sum.

\emph{Step 3 ($\cC$ off $\cN$: same points, two spacings).} $F_{t/5}-F_t=\sum_\ell(k_\ell^{t/5}-k_\ell^t)\prod_{\text{others}}$ (mixed products, every factor bounded by (L1) at its own spacing). Line differences at common points $x,u,v\in\Lambda_t$ with all distances $\ge8t$: translation-invariant factors $C,\nabla C,D,\nabla D$: $\le Ctw_{s+1}$ by (T2), slot $s\to s+1$; coarse factors $\Gamma,\Gamma^\ph$: $\le Ct$, $\nabla\Gamma,\nabla\Gamma^\ph$: $\le Ct(1+\ln t^{-1})^3$ by ($\Gamma$4), slots unchanged (for $f_D$, each member with its own weight, ($\Gamma$4) as stated); outer lines $G=C-\Gamma$, $\nabla G$: the sum of a translation-invariant and a coarse difference, hence $\le Ctw_{s+1}+Ct(1+\ln t^{-1})^3\le Ct(1+\ln t^{-1})^3(R+1)^{s+1}w_{s+1}$ on $|r|\le R$, and $\le2C\tau^4e^{-c|r|}$ on $|r|\ge R$ (both members are tails there, (O2), with the same $\alpha_0$ for both spacings). An argument shifted by one lattice unit at the two spacings ($te_\mu$ versus $(t/5)e_\mu$) is a displacement $\le t$ and falls under Step~4. For a differenced inner line the two outer lines still confine $u,v$ to $R$-balls (or are tails); for a differenced outer line its ``in'' part is confined by definition and its ``out'' part carries $e^{-c|r|}$: no mass decay is used. \cref{app:a04:lem:PC} with $\Sigma\le6$, $R\le C(1+\ln t^{-1})$, and the tail version: $|\cC\text{ off }\cN|\le Ct(1+\ln t^{-1})^{q}$.

\emph{Step 4 ($\cD$ off $\cN$: displaced arguments, one spacing $\tau=t/5$).} For $(x',u',v')\in c(x)\times c(u)\times c(v)$ every pairwise displacement is $\le2\sqrt3t\le|r|/4$ ($|r|\ge8t$), so $|r'|\ge|r|/2$ and (L1) at the displaced arguments is $\le2^s$(L1) at $(x,u,v)$; $R$-confinement holds with $R\pm4$. Telescope over the lines: translation-invariant factors by (T3): $\le Ctw_{s+1}(r)$; coarse factors $\Gamma,\Gamma^\ph$ by ($\Gamma$2) along lattice paths: $\le Ct$; $\nabla\Gamma,\nabla\Gamma^\ph$ by ($\Gamma$3): $\le Ct^\theta$ (slots unchanged; for $f_D$ the finest-scale modulation of $\nabla\Gamma$ is inside the constant of ($\Gamma$3)); outer lines: translation-invariant plus coarse parts as above, tails by (O2). \cref{app:a04:lem:PC} ($\Sigma\le6$): $|\cD\text{ off }\cN|\le C_\theta t^\theta(1+\ln t^{-1})^{q}$. Two-point terms: the same.

\emph{Step 5 ($b$).} $|\cT_{t/5}(\mathfrak a_{j+1})-\cT_{t/5}(\mathfrak a_j)|$: telescope over the lines; translation-invariant lines do not depend on $b$. Inner coarse factors ($\Gamma,\Gamma^\ph$, and $G_{\mathrm{in}}$ through $\Gamma$): ($\Gamma$5) gives $\le C|\mathfrak a_j-\mathfrak a_{j+1}|\le Ct^2$, confined by the two outer lines as in Steps~3--4. A differenced outer line: the same in/out dichotomy. \cref{app:a04:lem:CT} is uniform in $b\in[\mathfrak a(1-5^{-2}),\mathfrak a]$, so for $|r|\ge R$ both members obey (O2), $|k(\mathfrak a_{j+1})(r)-k(\mathfrak a_j)(r)|\le2C\tau^4e^{-c|r|}$; for $|r|\le R$ use ($\Gamma$5): $\le Ct^2(R+1)^sw_s(r)$, $s$ its slot. \cref{app:a04:lem:PC} and its tail version give $\le Ct^2(1+\ln t^{-1})^q$.

Summing Steps~2--5 and absorbing the logarithms into $t^\theta$ ($\theta<1$ arbitrary) gives (III-inc). Summation by parts (\cref{app:a04:lem:D}) was on the infinite lattice, so there is no boundary term, and all sums are absolutely convergent at fixed $\tau$ (outer lines decay exponentially by (O2), inner lines are bounded by $w_1$ at worst).
\end{proof}

\begin{proposition}[the sup bound at order $e^2$]\label{app:a04:prop:S}
For both weights, for all $k$ and $1\le j\le k$, with $C$ independent of $k$ and $j$ (depending on $m$, $\mathfrak a$, $m_0$, $\norm f_\infty$):
\[
\big\|e^2[\mathfrak g_x(C,C,C;D)-\mathfrak g_x(G_j,G_j,G_j;D_j)]\big\|_\infty=\norm{\beta_{(2,0)}+\mathfrak a}_\infty\le Ce^2\big[(1+(m\eta)^{-1})+(1+j\ln5)^{14}\big]\le Ce^2\eta^{-1}(1+j)^{14}.
\]
\end{proposition}
\begin{proof}
Pointwise (no cell average), at unit block scale with $\tau=t=5^{-j}$: $(\Gamma\Pi_CC)(x,x)=\int_{\mathbb T}\Mh\,\overline{\hat v_x}\,\hat u_x\dP$ with $\hat u_x\defeq\sum_{n\in N_\tau}s_\tau\hat C_\tau^2\Pih_\tau(p)e^{-ip\cdot x}$; by (E1)--(E3) and ($\Pi$-a), $|\hat u_x|\le C(|P|+m_0)(P^2+m'^2)^{-2}+C\sum_{n\ne0}\rho(n)|n|^{-3}\le C(|P|+m_0)(P^2+m'^2)^{-2}$ (for $f_D$ the D-aliases add $\sum_{\ell\ne0}\sigma(q)(25\tau^2/4)^2C\sqrt3\pi/\tau\le C\tau^3(1+\ln\tau^{-1})^3$), and $|\hat v_x|\Mh\le C$ ((A3) for $f_D$); so $|(\Gamma\Pi_CC)(x,x)|\le C\int_{\mathbb T}(|P|+m_0)(P^2+m'^2)^{-2}\dP\le C(1+\ln(1/m')+m_0/m')\le C(1+m'^{-1})$ (this is the infrared logarithm of the order-$e^2$ Wick constant). $(C\Pi_C\Gamma)(x,x)$: the same. $(\Gamma\Pi_C\Gamma)(x,x)=\int\Mh^2|\hat v_x|^2Y_\tau\dP\le C\int(P^2+m'^2)^2(P^2+m'^2)^{-2}(|P|+m_0)(P^2+m'^2)^{-2}\dP\le C(1+m'^{-1})$. $(\mathrm{III})$: (III-sup). By \eqref{app:a04:eq:dens}--\eqref{app:a04:eq:sandwich} the original-vertex density is $-2e^2[(\mathrm I)+(\mathrm{II})+(\mathrm{III})]$ at unit block scale; with $m'=m\eta$, $1+m'^{-1}=1+(m\eta)^{-1}$ and $(1+\ln t^{-1})^{14}=(1+j\ln5)^{14}$. The last inequality uses $\eta\le1$ and $(m\eta)^{-1}\le m^{-1}\eta^{-1}$.
\end{proof}

\begin{theorem}[the original-vertex part]\label{app:a04:thm:B2orig}
Let $\theta\in(0,1)$. There are $C$, $C_\theta$ (depending on $\theta$, $\mathfrak a$, $m_0\ge\max(m,\lmu)\eta$, $e$, $m$ and $\norm f_\infty$) such that for all $k$ and $1\le j\le k$ ($t=5^{-j}$, $\eta=5^j\eps_k\le1$), for both weights,
\begin{align*}
\norm{\beta_{(2,0)}+\mathfrak a}_\infty&\le Ce^2\big[(1+(m\eta)^{-1})+(1+j\ln5)^{14}\big]\le Ce^2\eta^{-1}(1+j)^{14},\\
\big|(\bar\beta_{(2,0)}+\bar{\mathfrak a})^{(j)}_k-(\bar\beta_{(2,0)}+\bar{\mathfrak a})^{(j+1)}_{k+1}\big|&\le C_\theta e^2\big[t(1+(m\eta)^{-1})+t^\theta\big]\le C_\theta e^2\,5^{-\theta j}\eta^{-1},
\end{align*}
where $\bar{\mathfrak a}$ is the block mean of $\mathfrak a_x$. This is the statement of \cref{app:a04:lem:B2} for the original-vertex part, with $\vartheta=\theta$ arbitrary in $(0,1)$ and $q=14$.
\end{theorem}
\begin{proof}
By \eqref{app:a04:eq:dens}--\eqref{app:a04:eq:sandwich}, $\beta_{(2,0)}+\mathfrak a=-2e^2[(\mathrm I)+(\mathrm{II})+(\mathrm{III})]$ at unit block scale, and the block mean $\eta^{-3}\sum_{x\in\Lk\cap[0,\eta)^3}\eps_k^3(\cdot)$ is the cell average $\langle\cdot\rangle_\cell$ of \cref{app:a04:lem:cell} at unit block scale. The diagonal pair $(k,j)\to(k+1,j+1)$ keeps $\eta,m',\mu'$ and the unit cell fixed and changes $(\tau,b,\text{weight})=(t,\mathfrak a_j,f\text{ on }5t\text{-blocks})\to(t/5,\mathfrak a_{j+1},f\text{ on }t\text{-blocks})$ (the reduction of \cref{app:a04:lem:B0} and its Dirichlet-weight paragraph). \Cref{app:a04:thm:R}: $|[(\mathrm I)+(\mathrm{II})]_{t/5}(\mathfrak a_{j+1})-[(\mathrm I)+(\mathrm{II})]_t(\mathfrak a_j)|\le Ct(1+m'^{-1})$. \Cref{app:a04:thm:III}: $|\cT_{t/5}(\mathfrak a_{j+1})-\cT_t(\mathfrak a_j)|\le C_\theta t^\theta$. Sup bound: \cref{app:a04:prop:S}. With $m'=m\eta$: $t(1+m'^{-1})\le(1+m^{-1})5^{-j}\eta^{-1}$ and $t^\theta=5^{-\theta j}\le5^{-\theta j}\eta^{-1}$. For $f_D$ the restriction $j\ge1$ ($\tau\le\frac15$ for both members) is the one in the statement of \cref{app:a04:lem:B2}.
\end{proof}

% a04_Bh_bulk4  Proof of app:a04:lem:B2, part 4: the averaging vertices, structure and line lemmas, the sup bound

\subsubsection*{B2.7 The averaging vertices: structure and the photon slices}

$\mathfrak a_x$ is the order-$e^2$ contribution to $\langle r_x^2\rangle^{\fl}$ of the $A$-dependence of King's block averaging, the vertices \cite[(3.24)--(3.26)]{K1} with the contour convention \cite[(2.19), (3.27)]{K1}. We work at unit block scale ($t=5^{-j}$, $\Lambda_t$, $m',\mu'\in(0,m_0]$, $b=\mathfrak a_j$) and write $\kappa_0\defeq\alpha_0/\sqrt3$, so that \cref{app:a04:lem:CT} at unit block scale reads $|G_\tau(x,y)|\le C\tau^{-3}e^{-\kappa_0|x-y|}$ and, at block side $\ell$, $|D_i(u,v)|\le C\ell^2\tau^{-3}e^{-\kappa_0|u-v|/\ell}$.

\emph{Photon slices.} For $0\le i\le j$ let $\ell_i\defeq5^it$, $P_i$ the uniform block average onto the blocks of side $\ell_i$ (corner blocks \cite[(2.12)]{K1}), $D_i\defeq(-\Delta_t+\mu'^2+\mathfrak a_i\ell_i^{-2}P_i)^{-1}$ ($D_0\defeq0$), and $S_i\defeq D_{i+1}-D_i$ for $i<j$. By \cite[(2.17)--(2.18)]{K1} the photon fluctuation field after $j$ steps is $A'=\sum_{i<j}A'_i$ with independent Gaussian slices, $\langle A'_{i,\mu}(u)A'_{i',\nu}(v)\rangle=\delta_{ii'}\delta_{\mu\nu}S_i(u,v)$ (\cref{app:a04:def:bulk}: $D_j=\sum_{i<j}G^\ph_{(i)}$).

\emph{Contours.} By \cite[(2.19)]{K1}, $A(\Gamma_{y,x})=\sum_{i<j}A'_i(\Gamma^{(i+1)}_{x_{i+1},x})$, $A(\Gamma)\defeq\sum_{\mathsf b\in\Gamma}tA_{\mathsf b}$, where $x_{i+1}$ is the corner of the $\ell_{i+1}$-block of $x$ and $\Gamma^{(i+1)}_{x_{i+1},x}$ is a staircase of $\le3\sum_{i'\le i}5^{i'+1}$ fine bonds, hence of length $\le C_L\ell_i$, $C_L\defeq3L^2/(L-1)$: slice $i$ sees only the sub-contour inside its own $\ell_{i+1}$-block.

\emph{Vertices.} With $U(\vartheta)=1+\vartheta J-\vartheta^2/2+O(\vartheta^3)$ ($J$ the complex structure on the two real components) and $V_{\mathrm{av}}=-\frac b2\sum_y[|Q(A)\phi|^2-|Q\phi|^2]$, $(Q(A)\phi)(y)=\sum_{x\in B(y)}t^3U(e'A(\Gamma_{y,x}))\phi(x)$ \cite[(3.24)--(3.26)]{K1} at zero block fields:
\begin{gather*}
V_1=-be'\sum_y\sum_{x\in B(y)}t^3A(\Gamma_{y,x})\,\ip{Q\phi(y)}{J\phi(x)},\qquad
V_2=\tfrac{be'^2}2\sum_{y,x}t^3A(\Gamma_{y,x})^2\ip{Q\phi(y)}{\phi(x)}-\tfrac{be'^2}2\sum_y\Big|\sum_{x\in B(y)}t^3A(\Gamma_{y,x})J\phi(x)\Big|^2,\\
V_h=e't^3\sum_{\mathsf b}A_{\mathsf b}\ip{\nabla\phi(\mathsf b)}{J\phi(\mathsf b_+)}=-e't^2\sum_{\mathsf b}A_{\mathsf b}\ip{\phi(\mathsf b_-)}{J\phi(\mathsf b_+)},
\end{gather*}
the last being the hopping vertex \cite[(3.22)]{K1} at $m=1$ (one difference quotient) at unit block scale ($\ip{\phi(\mathsf b_+)}{J\phi(\mathsf b_+)}=0$; its sign plays no role). At order $e'^2$ of $\langle|\phi(x_0)|^2\rangle^{\fl}$, since $\langle A'\rangle=0$ kills the first order and everything without an averaging vertex is the original-vertex part, $\mathfrak a_{x_0}$ consists of the four classes of connected correlations
\[
\mathrm{E1}+\mathrm{E2}=\langle|\phi(x_0)|^2;V_2\rangle,\qquad\mathrm{E3}=\tfrac12\langle|\phi(x_0)|^2;V_1;V_1\rangle,\qquad\mathrm{E4}=\langle|\phi(x_0)|^2;V_1;V_h\rangle.
\]
Their photon kernels are $\cV(x)\defeq\sum_i\Var(A'_i(\Gamma^{(i+1)}_{x_{i+1},x}))$, $\cK(x,x')\defeq\sum_i\langle A'_i(\Gamma(x))A'_i(\Gamma(x'))\rangle$ and $\cK_{\mathsf b}(x)\defeq\sum_i\cK_{i,\mathsf b}(x)$, $\cK_{i,\mathsf b}(x)\defeq\sum_{\mathsf b'\in\Gamma_i(x),\ \mathsf b'\parallel\mathsf b}tS_i(\mathsf b',\mathsf b)$; the scalar lines are $G=G_t$, $(QG)$, $(QGQ^*)$. Each class is a two- or three-line cycle through $x_0$ (E1/E2: two lines from $x_0$ to legs in one block; E3/E4: $x_0$--$v_1$--$v_2$--$x_0$), and E4 carries one differentiated line: the $J$-antisymmetry at $\mathsf b$ gives $2(QG)(y,x_0)[G(x_0,\mathsf b_-)G(x,\mathsf b_+)-G(x_0,\mathsf b_+)G(x,\mathsf b_-)]$, a forward difference at $\mathsf b$ on one of the two lines.

\begin{lemma}[photon slices]\label{app:a04:lem:S}
\begin{enumerate}[label=\textup{(S\arabic*)},leftmargin=*,start=0]
\item $S_i$ is the slice covariance of \cite[(2.17)]{K1}: the joint Gaussian of $(\phi,\psi_i)$ there has marginal $\phi$-covariance $D_{i+1}$, because $\ell_{i+1}^2/\mathfrak a_{i+1}=\ell_i^2/\mathfrak a_i+\ell_{i+1}^2/\mathfrak a$ (algebra from \cite[(2.13)]{K1}), and by the law of total covariance $D_{i+1}=\Cov(\phi|\psi)+\Cov(\E[\phi|\psi])=D_i+\mathfrak a_i^2\ell_i^{-4}D_iQ_i^*C^{(i)}Q_iD_i$, the term of \cite[(2.17)]{K1}.
\item $0\le S_i\le D-D_i$: $P_{i+1}=Q_i^*Q'^*Q'Q_i\le P_i$ ($Q'Q'^*=1$) and $\mathfrak a_i\ell_i^{-2}\ge\mathfrak a_{i+1}\ell_{i+1}^{-2}$, so $D_i\le D_{i+1}\le D$.
\item $|S_i(u,v)|\le C\ell_i^{-1}$: $S_i(u,u)\le(D-D_i)(u,u)$ by (S1), Cauchy--Schwarz, and $(D-D_i)(u,u)=\ell_i^{-1}\Gamma_{t/\ell_i}(u/\ell_i,u/\ell_i)\le C\ell_i^{-1}$ by \cref{app:a04:lem:B0}(i) at relative spacing $t/\ell_i=5^{-i}$, mass $\mu'\ell_i\le m_0$, coupling $\mathfrak a_i$; for $i=0$, $D(u,u)\le C/t$. Uniform in $\mu'$.
\item $|\nabla_uS_i(u,v)|\le C\ell_i^{-2}$: ($\Gamma$2) rescaled for $D-D_i$ and $D-D_{i+1}$ ($i\ge1$), (T1) for $i=0$.
\item $|S_i(u,v)|\le C\ell_i^{-1}\min\big(1,e^{(\kappa_0/2)(\rho_i-|u-v|)/\ell_{i+1}}\big)$ with $\rho_i\defeq(6/\kappa_0)\ell_{i+1}\ln(\ell_{i+1}/t)\le C\ell_i(i+1)$: \cref{app:a04:lem:CT} for $D_i$ and $D_{i+1}$ at block side $\le\ell_{i+1}$, half the exponent paying the prefactor $\ell_{i+1}^2t^{-3}$ beyond $\rho_i$.
\item For $u,v\in\Lambda_t\subset\Lambda_{t/5}$ and $i\ge1$: $|S^{(t)}_i(u,v)-S^{(t/5)}_{i+1}(u,v)|\le Ct\ell_i^{-2}$: the coarse slice $i$ and the fine slice $i+1$ have the same block sides $\ell_i,\ell_{i+1}$; compare each of the four coarse lines with the continuum one at its block side by \cref{app:a04:lem:B0}(ii) rescaled ($\ell^{-1}\cdot C(t/\ell)$), and the coupling change $\mathfrak a_{i}\to\mathfrak a_{i+1}$ by \cref{app:a04:lem:B0}(iv) rescaled ($C\ell_i^{-1}5^{-2i}=Ct^2\ell_i^{-3}\le Ct\ell_i^{-2}$). Slice $0$ of both members and the extra finest slice of the finer member are not matched; each contributes $O(t\cdot\mathrm{polylog})$ (E4: $O(t^2\cdot\mathrm{polylog})$) by \cref{app:a04:lem:C}(i),(ii) below, since $\ell_0=t$ for both.
\end{enumerate}
\end{lemma}

\begin{lemma}[contour kernels]\label{app:a04:lem:C}
\begin{enumerate}[label=\textup{(\roman*)},leftmargin=*]
\item $\Var(A'_i(\Gamma^{(i+1)}))\le(C_L\ell_i)^2C\ell_i^{-1}=C\ell_i$; $\cV\le C\sum_{i<j}\ell_i\le C$; $|\cK_i|\le C\ell_i$ and $|\cK|\le C$ (Cauchy--Schwarz).
\item $|\cK_{i,\mathsf b}(x)|\le C$ ($\le C_L\ell_i/t$ bonds of weight $t$, (S2)); up to the tail of (S4) it is supported in $|\mathsf b-x|\le r_i\defeq C_L\ell_i+\rho_i\le C\ell_i(i+1)$.
\item For $x\in\Lambda_t$ (corner blocks, $x_1=x$) the finer member's sub-contours coincide with the coarse ones as point sets and its finest slice sees the empty contour; for $x'$ in the coarse cell of $x$ every slice gains $\Gamma_{x,x'}$ ($\le3(L-1)$ fine bonds, length $\le3t$).
\item Replacing the fine bond sum along a contour by the coarse one in the same kernel costs $\le2(C_L\ell_i)^2C\ell_i^{-2}t=Ct$ per slice ((S3)).
\item In-cell oscillation of the finer member's kernels: $|\cK'(x',\cdot)-\cK'(x,\cdot)|\le Ct(j+1)$; for $\cK'_{\mathsf b}$ the E4 cost is $\le Ctj^2$.
\end{enumerate}
\end{lemma}
\begin{proof}
(i): the variance of $A'_i$ along a contour of $\le C_L\ell_i/t$ bonds of length $t$ is bounded by (length)$^2\sup|S_i|$, (S2); $\sum_{i<j}\ell_i=t(5^j-1)/4\le\frac14$. (ii): (S2) and (S4), the bonds of $\Gamma_i(x)$ lying within $C_L\ell_i$ of $x$. (iii) is the corner-block convention \cite[(2.12), (2.19)]{K1}. (iv): (S3), the Riemann error of a bond sum with step $t$ over a path of length $C_L\ell_i$ against a kernel with gradient $\le C\ell_i^{-2}$. (v): (iii) and (S2), (S3), summed over the $\le j+1$ slices.
\end{proof}

\begin{lemma}[scalar lines]\label{app:a04:lem:G}
Write $|r|_t\defeq|r|+t$ and $R_t\defeq(6/\kappa_0)\ln t^{-1}=6\sqrt3\,\alpha_0^{-1}\ln t^{-1}$, the radius used in the proof of \cref{app:a04:lem:B1} (in \S\S B2.7--B2.10 $R_t$ denotes this radius, not the $R_\tau$ of (O2) in \cref{app:a04:lem:TIOUT}); $G=C_t-\Gamma_t$.
\begin{enumerate}[label=\textup{(G\arabic*)},leftmargin=*]
\item $|C_t(r)|\le C|r|_t^{-1}$, $|\nabla C_t|\le C|r|_t^{-2}$, $|\nabla\nabla C_t|\le C|r|_t^{-3}$ (T1); $|C_t-C_{t/5}|(r)\le Ct^2|r|^{-3}$ and $|\nabla^{(t)}C_t-\nabla^{(t/5)}C_{t/5}|(r)\le Ct|r|^{-3}$ for $r\in\Lambda_t\setminus0$ (\cref{app:a04:lem:LP,app:a04:lem:LPgrad} at both spacings against $C_0$, $\partial C_0$).
\item $|\Gamma_t|,|\nabla\Gamma_t|\le C$; $|\Gamma_t-\Gamma_{t/5}|\le Ct+C5^{-2j}$ (($\Gamma$4) and ($\Gamma$5) for the change of $b$); $|\partial_b\Gamma_t|\le C$; $|\nabla\nabla\Gamma_\tau(u,v)|\le C\sum_{n\in N_\tau}|p|^2\sigma\hat C_\tau\le C\sum_{|n|\le C/\tau}\rho(n)\le C(1+\ln\tau^{-1})^3$ (uniform in $m'$).
\item $|G(x,y)|\le Ct^{-3}e^{-\kappa_0|x-y|}$ (\cref{app:a04:lem:CT}); hence $|G|\le C|r|_t^{-1}+C$ on $|r|\le R_t$ and $\le Ct^3$ beyond; $|(QG)|,|(QGQ^*)|\le C$ with the same tail; $t^3G(x,x)\le Ct^2$.
\item Gradient line: $|\nabla G|\le C|r|_t^{-2}+C$; $|\nabla_{\mathsf b}(QG)(y,\mathsf b)|\le C$ uniformly in $y$.
\end{enumerate}
\end{lemma}
\begin{proof}
(G1), (G2) restate \cref{app:a04:lem:TIOUT,app:a04:lem:Gam,app:a04:lem:LP,app:a04:lem:LPgrad}; the second-difference bound on $\Gamma_\tau$ follows from the Fourier representation with two factors $|d_\tau(p_\mu)|\le|p|$ and (E1)--(E3). (G3), (G4): (G1), (G2) and \cref{app:a04:lem:CT}; $(QG)(y,\cdot)$ is a block average of $G$.
\end{proof}

\begin{lemma}[tails of the contour kernels]\label{app:a04:lem:T}
Fix $i\ge1$, $r\defeq|\mathsf b-x|$, $\lambda_i\defeq2\ell_{i+1}/\kappa_0$, $M\defeq(2/\kappa_0)\ln t^{-1}$ and $\varrho_i\defeq r_i+M\ell_{i+1}$; for $x,\mathsf b$ coarse let $\Delta\cK_{i,\mathsf b}(x)\defeq\cK^{(t)}_{i,\mathsf b}(x)-\cK^{(t/5)}_{i+1,\mathsf b}(x)$ (coarse slice $i$ matched with fine slice $i+1$).
\begin{enumerate}[label=\textup{(T\arabic*)},leftmargin=*]
\item $|\Delta\cK_{i,\mathsf b}(x)|\le Ct/\ell_i$ for every $r$ ((S5) and \cref{app:a04:lem:C}(iv): $\le(C_L\ell_i/t)\cdot t\cdot Ct\ell_i^{-2}$).
\item $|\cK^{(t)}_{i,\mathsf b}(x)|\le Ce^{-(r-r_i)/\lambda_i}$ and $|\cK^{(t/5)}_{i+1,\mathsf b}(x)|\le CL^3e^{-(r-r_i)/\lambda_i}$ for $r>r_i$, hence $|\Delta\cK_{i,\mathsf b}(x)|\le CL^3e^{-(r-r_i)/\lambda_i}$.
\item (E4 cost, both Wick families.) Let $\kappa\ge0$ be radial and non-increasing with $\kappa(r)\ge|\Delta\cK_{i,\mathsf b}(x)|\mathbf 1_{r>r_i}$, e.g.\ $\kappa_i(r)\defeq\min(Ct/\ell_i,CL^3e^{-(r-r_i)_+/\lambda_i})$. Then the E4 term with the kernel $\Delta\cK_{i,\mathsf b}\mathbf 1_{r>r_i}$ is $\le Ct(1+j)^3R_t^3$ per slice, and the same holds for the kernel $|\cK^{(\tau)}_{i,\mathsf b}|\mathbf 1_{r>r_i}$ of a single member with $Ct$ replaced by $C\ell_{i+1}$.
\end{enumerate}
\end{lemma}
\begin{proof}
(T2): bound each member separately. Coarse member: the contour $\Gamma^{(i+1)}_{x_{i+1},x}$ ends at $x$ and has length $\le C_L\ell_i$, so each of its bonds $\mathsf b'$ has $|\mathsf b'-\mathsf b|\ge r-C_L\ell_i$, and (S4) gives $|S_i(\mathsf b',\mathsf b)|\le C\ell_i^{-1}e^{(\kappa_0/2)(\rho_i+C_L\ell_i-r)/\ell_{i+1}}=C\ell_i^{-1}e^{-(r-r_i)/\lambda_i}$; $\le C_L\ell_i/t$ bonds of weight $t$ give $\le Ce^{-(r-r_i)/\lambda_i}$. Finer member (slice $i+1$ of $\Lambda_{t/5}$): same scale $\ell_i$, same block side $\ell_{i+1}$, contour length $3(t/5)\sum_{i'\le i+1}5^{i'+1}\le C_L\ell_i$; its (S4) shift is $(6/\kappa_0)\ell_{i+1}\ln(\ell_{i+1}/(t/5))=\rho_i+(6/\kappa_0)\ell_{i+1}\ln5$, which multiplies the bound by $e^{3\ln5}=L^3$.
(T3): E4 contracts $|\phi(x_0)|^2$ against $V_1\ni\ip{Q\phi(y)}{J\phi(x)}$ and $V_h\ni\ip{\phi(\mathsf b_-)}{J\phi(\mathsf b_+)}$ (self-loops vanish, $\Tr J=0$; $x_0$ joined to both fields of one vertex leaves a self-loop on the other), so there are two Wick families, both carrying the kernel $\cK_{i,\mathsf b}(x)$: (i) $x_0$--$Q\phi(y)$, $x$--$\mathsf b$, and (ii) $x_0$--$\phi(x)$, $Q\phi(y)$--$\mathsf b$. \emph{Family (i).} By (G3)--(G4), uniformly in $\mathsf b$, for both placements of the difference (undifferentiated or gradient line from $\mathsf b$ to $x$), $\sum_xt^3(|G(x,\mathsf b_+)|+|\nabla G(x,\mathsf b)|)\kappa(|x-\mathsf b|)\le C\int(r^{-2}+r^{-1}+1)\kappa(r)d^3r+Ct\sup\kappa$ (the $O(1)$ points with $|x-\mathsf b|\le2t$ cost $\le t^3Ct^{-2}$ each). The outer factors: $\sum_{\mathsf b}t^3(|\nabla G(\mathsf b,x_0)|+|G(x_0,\mathsf b)|)\le CR_t^3$ (beyond $R_t$ a difference quotient is $\le2t^{-1}Ct^{-3}e^{-\kappa_0(|r|-t)}$, total $\le Ct^2R_t^2$) and $|(QG)(y,x_0)|\le C$. With $r_i\le C\ell_i(i+1)$, $M\ell_{i+1}=(2L\ln L/\kappa_0)\ell_ij$ and $\ell_i\le1$: $\varrho_i\le C\ell_i(1+j)$ and $\lambda_i\le2L/\kappa_0$. Shell $r_i<r\le\varrho_i$: $\kappa\le Ct/\ell_i$ by (T1), $\int_{r\le\varrho}(r^{-2}+r^{-1}+1)d^3r\le C\varrho(1+\varrho)^2$, cost $\le C(t/\ell_i)\ell_i(1+j)(1+j)^2=Ct(1+j)^3$. Region $r>\varrho_i$: by (T2), $\kappa\le CL^3e^{-(\kappa_0/2)M}e^{-(r-\varrho_i)/\lambda_i}=CL^3t\,e^{-(r-\varrho_i)/\lambda_i}$, and $\int_{r>\varrho}(r^{-2}+r^{-1}+1)e^{-(r-\varrho)/\lambda}d^3r\le C\lambda(1+\varrho+\lambda)^2\le C(1+j)^2$: cost $\le Ct(1+j)^2$. Near-diagonal: $Ct\sup\kappa\le Ct$. \emph{Family (ii)} (lines $G(x,x_0)$, $G(x_0,\mathsf b_\mp)$, $(QG)(y,\mathsf b_\pm)$, $y$ the block of $x$): the $J$-antisymmetry at $\mathsf b$ puts one forward difference at $\mathsf b$ on $G(x_0,\cdot)$ or on $(QG)(y,\cdot)$, and $|(QG)(y,\mathsf b)|,|\nabla_{\mathsf b}(QG)(y,\mathsf b)|\le C$ (G3)--(G4). The line from $x$ now runs to $x_0$. Split the $x$-sum: on $|x-x_0|\ge|x-\mathsf b|$, $|G(x,x_0)|\le C|x-\mathsf b|_t^{-1}+C$ ($|\cdot|_t^{-1}$ non-increasing); on $|x-x_0|<|x-\mathsf b|$, $\kappa(|x-\mathsf b|)\le\kappa(|x-x_0|)$. Hence, uniformly in $\mathsf b$ and $x_0$, $\sum_xt^3|G(x,x_0)|\kappa(|x-\mathsf b|)\le2\sup_z\sum_xt^3(C|x-z|_t^{-1}+C)\kappa(|x-z|)\le C\int(r^{-1}+1)\kappa(r)d^3r+Ct\sup\kappa$, which for $\kappa=\kappa_i$ is dominated by the three regions above ($r^{-1}+1\le r^{-2}+r^{-1}+1$; the first region's integral already runs over the whole ball $r\le\varrho_i$). The outer factor $\sum_{\mathsf b}t^3(|\nabla G(\mathsf b,x_0)|+|G(x_0,\mathsf b)|)\le CR_t^3$ is family (i)'s. Same per-slice cost $Ct(1+j)^3R_t^3$. For a single member, (T2) alone gives $\kappa\le Ce^{-(r-r_i)/\lambda_i}$ on $r>r_i$ and $\int_{r>r_i}(r^{-2}+r^{-1}+1)e^{-(r-r_i)/\lambda_i}d^3r\le C\lambda_i(1+r_i+\lambda_i)^2\le C\ell_{i+1}(1+j)^2$. The same holds for any radial non-increasing $\kappa$.
\end{proof}

\subsubsection*{B2.8 The sup bound for the averaging part}

\begin{proposition}[the averaging part, sup bound]\label{app:a04:prop:AVa}
For $f=1$, all $k$ and $1\le j\le k-J_0$, there are $q$ and $C$ (uniform in $m',\mu'\in(0,m_0]$, $b$, $k$, $j$) with $\sup_x|\mathfrak a_x|\le Ce^2(1+\ln t^{-1})^q$ at unit block scale, i.e.\ $\norm{\mathfrak a}_\infty\le Ce^2\eta^{-1}(1+j)^q$.
\end{proposition}
\begin{proof}
Bound each class by $|\text{photon kernel}|\times|\text{lines}|$, with the vertex sums restricted to $|\cdot-x_0|\le R_t$; the tails cost $\le Ct^3R_t^q$ by (G3). \emph{E1:} $2be'^2\sum_xt^3\cV(x)(QG)(y,x_0)G(x,x_0)\le C\sum_{|x-x_0|\le R_t}t^3(|x-x_0|_t^{-1}+1)\le C(1+R_t)^3$ ($\cV\le C$, \cref{app:a04:lem:C}(i), $|(QG)|\le C$). \emph{E2, E3:} $|\cK|\le C$ and $\sum t^3t^3(|x_0-u|_t^{-1}+1)(|u-v|_t^{-1}+1)(|v-x_0|_t^{-1}+1)$ over the $R_t$-balls $\le C(1+R_t)^6$ (local integrability in three dimensions; coincident points cost $Ct^2$). \emph{E4:} per slice $i$, with \cref{app:a04:lem:C}(ii) in its support form, $\sum_{\mathsf b}t^3|\nabla G(\mathsf b,x_0)|\sum_{|x-\mathsf b|\le r_i}t^3|G(x,\mathsf b_+)|\,|(QG)(y,x_0)|\le Cr_i^2(1+r_i)(1+R_t)^3$ (the ball sum of $|G(x,\mathsf b_+)|\le C|x-\mathsf b|_t^{-1}+C$ over $|x-\mathsf b|\le r_i$ is $\le C(r_i^2+r_i^3)$), the tail $|\mathsf b-x|>r_i$ costs $\le C\ell_{i+1}(1+j)^2R_t^3$ (\cref{app:a04:lem:T}(T3), single member), summable over $i$, and the other placements of the difference and the second Wick family are identical ($|\nabla_{\mathsf b}(QG)(y,\mathsf b)|\le C$). Summing over $i<j$, $\sum_ir_i^2(1+r_i)\le C(1+j)^3$ (dominated by $i\approx j$). Hence $\sup_x|\mathfrak a_x|\le Ce^2(1+\ln t^{-1})^q$ at unit block scale, with $e'^2=e^2\eta$ and the factor $\eta^{-1}$ of the density; $q$ is not pinned, and only its existence is used.
\end{proof}

% a04_Bh_bulk5  Proof of app:a04:lem:B2, part 5: the averaging part along the diagonal pairs, the Dirichlet weight, conclusion

\subsubsection*{B2.9 The averaging part along the diagonal pairs}

\begin{lemma}[gradient of the coarse-line pair difference]\label{app:a04:lem:Lam}
Let $g\defeq(\Gamma_t-\Gamma_{t/5})|_{\Lambda_t\times\Lambda_t}$ (both lines with the same $b$). Then
\[
|\nabla^{(t)}\Gamma_t-\nabla^{(t)}\Gamma_{t/5}|\le Ct^{1/2}(1+\ln t^{-1})^3,\qquad|\nabla^{(t)}\Gamma_{t/5}-\nabla^{(t/5)}\Gamma_{t/5}|\le Ct(1+\ln t^{-1})^3\quad\text{on }\Lambda_t\times\Lambda_t.
\]
\end{lemma}
\begin{proof}
$A\defeq\sup|g|\le Ct$ by ($\Gamma$4), and $B\defeq\sup|\nabla^{(t)}\nabla^{(t)}g|\le C(1+\ln t^{-1})^3$ by (G2) (coarse second differences of $\Gamma_{t/5}$ are averages of fine ones). Discrete Landau inequality: $g(u+Nte_\mu)-g(u)=Nt\,\nabla^{(t)}_\mu g(u)+t^2\sum_{i<N}\sum_{i'<i}\nabla\nabla g$, so $|\nabla^{(t)}g|\le2A/(Nt)+NtB/2$; take $N=\lfloor t^{-1/2}\rfloor$. The second bound: a coarse difference quotient of $\Gamma_{t/5}$ is the average of five fine ones along the step, and consecutive fine ones differ by $\le(t/5)B$.
\end{proof}

\begin{proposition}[the averaging part, increment]\label{app:a04:prop:AVb}
For $f=1$, all $k$ and $1\le j\le k-J_0$, with $q$ as in \cref{app:a04:prop:AVa} (enlarged if necessary) and $C$ uniform in $m',\mu'\in(0,m_0]$:
\[
\big|\bar{\mathfrak a}^{(j)}_k-\bar{\mathfrak a}^{(j+1)}_{k+1}\big|\le Ce^2\eta^{-1}\big[t^{1/2}(1+j)^q+5^{-2j}\big].
\]
\end{proposition}
\begin{proof}
Both members are evaluated at unit block scale on the same unit cell; the coarse member has spacing $t$, coupling $\mathfrak a_j$ and slices $i<j$, the finer one spacing $t/5$, coupling $\mathfrak a_{j+1}$ and slices $i'\le j$. Every term keeps a cycle through $x_0$ with unit-scale-decaying lines, so all vertex sums converge uniformly in $m',\mu'$ (the mass-free decay (G3)). Telescope one factor at a time in four steps.

\emph{(1) Unmatched slices} (slice $0$ of the coarser member, the two finest slices of the finer member): $O(t\cdot\mathrm{polylog})$ by (S5) and \cref{app:a04:lem:C}(i),(ii) ($\ell_0=t$ in both members).

\emph{(2) Common slices $i\ge1$ at coarse vertex points.} (S5) and \cref{app:a04:lem:C}(iv) give $Ct$ per slice on $\cV$ and $\cK$ (total $Ctj$). On $\cK_{\mathsf b}$: per slice $|\Delta\cK_{i,\mathsf b}|\le Ct/\ell_i$ (\cref{app:a04:lem:T}(T1)), which is not small for small $i$; in the region $|\mathsf b-x|\le r_i$ the E4 cost is $(t/\ell_i)\,r_i^2(1+r_i)R_t^3$ (the ball sum of $|G(x,\mathsf b_+)|\le C|x-\mathsf b|_t^{-1}+C$ over $|x-\mathsf b|\le r_i$ is $\le C(r_i^2+r_i^3)$), and $\sum_i(t/\ell_i)r_i^2(1+r_i)R_t^3\le Ct\sum_i\ell_i(i+1)^3R_t^3\le Ct(1+j)^3R_t^3$ ($r_i\le C\ell_i(i+1)$); in the shell $|\mathsf b-x|>r_i$ the pointwise bound carries no decay and would only give $\sum_it/\ell_i=O(1)$, so there \cref{app:a04:lem:T}(T3) is used, with the per-slice cost $Ct(1+j)^3R_t^3$ for both Wick families: the shell part is $\le Ct(1+j)^4R_t^3$.

\emph{(3) Lines at coarse points.} $C$-part: $|C_t-C_{t/5}|\le Ct^2|r|^{-3}$ and $|\nabla^{(t)}C_t-\nabla^{(t/5)}C_{t/5}|\le Ct|r|^{-3}$ (G1), summable against one more line at the cost of one logarithm; $\Gamma$-part: $|\Gamma_t-\Gamma_{t/5}|\le Ct+C5^{-2j}$ (G2); gradient-$\Gamma$ part: \cref{app:a04:lem:Lam}, $\le Ct^{1/2}(1+\ln t^{-1})^3$. This is the only step with rate $t^{1/2}$.

\emph{(4) Fine vertex points $\to$ coarse.} Kernel oscillation by \cref{app:a04:lem:C}(v) ($\le Ctj^2$ for $\cK'_{\mathsf b}$ in the region $|\mathsf b-x|\le r_i$); line oscillation $t(|r|_t^{-2}+1)$; gradient-line oscillation $Ct|r|_t^{-3}+Ct(1+\ln t^{-1})^3$ (G1), (G2); near-diagonal points ($|r|<2t$) bounded directly, $O(t)$. Shell: by \cref{app:a04:lem:C}(iii) the in-cell difference $D_{i'}(x',\mathsf b)\defeq\cK'_{i',\mathsf b}(x')-\cK'_{i',\mathsf b}(x)$ ($x'$ in the coarse cell of $x$) is the sum of $(t/5)S'_{i'}(\mathsf b',\mathsf b)$ over the $\le3(L-1)$ fine bonds $\mathsf b'$ of $\Gamma_{x,x'}$, all within $3t$ of $x$. With $\ell'\defeq5^{i'}t/5$ the scale of fine slice $i'$ ($0\le i'\le j$), $\lambda'\defeq2L\ell'/\kappa_0$, $\rho'\defeq(6/\kappa_0)L\ell'\ln(L^2\ell'/t)\le C\ell'(1+j)$ and $r'\defeq\rho'+3t$, (S2) and (S4) of the finer member give $|D_{i'}(x',\mathsf b)|\le C(t/\ell')\min(1,e^{-(r-r')/\lambda'})$, $r=|\mathsf b-x|$. The bond move ($\mathsf b\to$ the parallel bond based at the coarse point of its cell, $|\mathsf b-\mathsf b_c|\le3t$) obeys the same form $\min(Ct/\ell',CL^3e^{-(r-r')/\lambda'})$: pointwise by (S3) ($|\nabla S'_{i'}|\le C\ell'^{-2}$, $\le C_L\ell'L/t$ bonds of weight $t/5$, step $\le3t$), tail by (S4) applied to each of the two terms. \cref{app:a04:lem:T}(T3) for both Wick families ($t/\ell'\le L$; for the bond move with the split at $r'+ML\ell'$) gives per slice $\le Ct(1+j)^3$, over the $\le j+1$ fine slices $\le Ct(1+j)^4$, times the outer factor $CR_t^3$.

\emph{Coupling change.} $|\mathfrak a_j-\mathfrak a_{j+1}|\le C5^{-2j}$ with (G2) ($|\partial_b\Gamma|\le C$) and (S5): $C5^{-2j}$.

Summing (1)--(4): $|\bar{\mathfrak a}^{(j)}_k-\bar{\mathfrak a}^{(j+1)}_{k+1}|\le Ce^2\eta^{-1}[t^{1/2}(1+j)^q+5^{-2j}]$, the factors $e^2\eta^{-1}$ as in \cref{app:a04:prop:AVa}.
\end{proof}

\subsubsection*{B2.10 The averaging part with the Dirichlet weight, and the conclusion}

\begin{proposition}[the averaging part, weight $f_D$]\label{app:a04:prop:AVf}
\Cref{app:a04:prop:AVa,app:a04:prop:AVb} hold for the weight $f_D$ of \cref{app:a04:def:bulk}, with constants depending in addition on $\norm{f_D}_\infty$.
\end{proposition}
\begin{proof}
\emph{Setting.} The photon kernels are unchanged (the photon averaging is unchanged, \cite[p.~328]{K2}, so \cite[(2.17)]{K1} and \cref{app:a04:lem:S} hold verbatim). The vertex is $V^f_{\mathrm{av}}=-\frac b2\sum_y[|Q^f(A)\phi|^2-|Q^f\phi|^2]$ with $(Q^f(A)\phi)(y)=\sum_{x\in B(y)}t^3f(x)U(e'A(\Gamma_{y,x}))\phi(x)$ \cite[(3.2)]{K2}, the weight multiplying the fine field at every slice; $Q^f_{j+1}=Q_1Q^f_j$, $P^f_{i+1}\le P^f_i$ as before; the weight contributes bounded factors $\le\norm{f_D}_\infty^2$. \emph{Lines.} $G^f=C_t-\Gamma^f_t$ with the Dirichlet-weight clause of \cref{app:a04:lem:B0} (the same $\Gamma_0$, lower bound (A3) for $j\ge1$), \cref{app:a04:lem:CT} with $f_D$, and the gradient and second-difference bounds of \cref{app:a04:lem:G} with the D-aliases contributing $\le|p|^i\,Ct^2(1+\ln t^{-1})^3\le C\cdot\mathrm{polylog}$ for $i\le2$ ($|p|\le C/t$, $\hat C_t\le25t^2/4$, (A1)). With these inputs the proofs of \cref{app:a04:prop:AVa,app:a04:prop:AVb} go through verbatim, except for the weight mismatch of the pair: $f^{(t)}$ is $5t$-periodic with $5t$-block average $1$, $f^{(t/5)}$ is $t$-periodic with $t$-cell average $1$ \cite[(3.1)]{K2}. \emph{Weight normalisation.} For an integrand $F'$ on the finer lattice and $F$ on the coarser one,
\begin{gather*}
\textstyle\sum_{x'}(t/5)^3f'(x')F'(x')=\sum_{x\in\Lambda_t}t^3F'(x)+\sum_xt^35^{-3}\sum_{x'\in c(x)}f'(x')[F'(x')-F'(x)],\\
\textstyle\sum_xt^3f(x)F(x)=\sum_xt^3F(x)+\sum_{\text{blocks }B}\sum_{x\in B}t^3(f(x)-1)[F(x)-F(x_B)].
\end{gather*}
Both remainders are $\le\norm f_\infty$ times the in-cell, resp.\ in-block, oscillation of the integrand, which is step (4) of the proof of \cref{app:a04:prop:AVb} over $5t$-blocks: $O(t\cdot\mathrm{polylog})$ (near-diagonal points $\le Ct^2$). The same reduction applies to the $(Q^fG)$ lines. Hence the increment bound holds for $f_D$ with the same rate $t^{1/2}$.
\end{proof}

\begin{proof}[Proof of \cref{app:a04:lem:B2}]
By \cref{app:a04:lem:bulkform}, $\beta_{(2,0)}=(\beta_{(2,0)}+\mathfrak a)-\mathfrak a$, the first bracket being the original-vertex part. For it, \cref{app:a04:thm:B2orig} gives the sup bound $Ce^2\eta^{-1}(1+j)^{14}$ and the increment bound $C_\vartheta e^25^{-\vartheta j}\eta^{-1}$ for every $\vartheta\in(0,1)$, for both weights, for $1\le j\le k$. For the averaging part, \cref{app:a04:prop:AVa,app:a04:prop:AVb,app:a04:prop:AVf} give $\norm{\mathfrak a}_\infty\le Ce^2\eta^{-1}(1+j)^q$ and $|\bar{\mathfrak a}^{(j)}_k-\bar{\mathfrak a}^{(j+1)}_{k+1}|\le Ce^2\eta^{-1}[5^{-j/2}(1+j)^q+5^{-2j}]\le C_\vartheta e^25^{-\vartheta j}\eta^{-1}$ for every $\vartheta\in(0,\frac12)$, for both weights, for $1\le j\le k-J_0$. Adding, with $q$ the larger of $14$ and the exponent of the averaging part, gives both bounds of \cref{app:a04:lem:B2} with the constants depending on $\vartheta$, $\mathfrak a$, $e$, $m$, $\lmu$ and $m_0$ (and on $\norm{f_D}_\infty\le(5/3)^3$).
\end{proof}


\subsubsection{The bulk sequence}

\begin{lemma}[the bulk sequence]\label{app:a04:lem:BULK}
Fix $\gamma_b\in(0,\frac12)$ and put $\beta_b\defeq1$. There is $C$ such that for all $k$ and $1\le j\le k-J_0$,
\begin{equation}\label{app:a04:eq:BULKinc}
\big|\bar\beta^{(j)}_k-\bar\beta^{(j+1)}_{k+1}\big|\ \le\ C\,5^{-\gamma_bj}\eta^{-\beta_b},
\end{equation}
and along the diagonal pairs, for $k\ge k_0$ and $j=j(k)$,
\begin{equation}\label{app:a04:eq:BULKsup}
\norm{\beta^{(j)}_k}_\infty\ \le\ C\,\eta_k^{-2}.
\end{equation}
\end{lemma}
\begin{proof}
Add the bounds of \cref{app:a04:lem:B0,app:a04:lem:B1,app:a04:lem:B2} with $\gamma''=\vartheta=\gamma_b$: the increments are $\le C\eta^{-1}(5^{-j/2}+5^{-\gamma_bj}+5^{-2j})\le C5^{-\gamma_bj}\eta^{-1}$ (as $\eta\le1$). For the sup bound along the diagonal pairs, $j(k)\le k$ and $\eta_k^{-1}=5^{k-j(k)}\ge5^{\alpha k-1}$, so $(1+j)^3\le(1+k)^3\le C_\alpha\eta_k^{-1}$ and $(1+j)^q\le(1+k)^q\le C_\alpha\eta_k^{-1}$, and $\norm{\beta^{(j)}_k}_\infty\le C\eta_k^{-1}+C\lambda(1+k)^3+Ce^2\eta_k^{-1}(1+k)^q\le C\eta_k^{-2}$.
\end{proof}

\subsection{Proof of \texorpdfstring{\cref{prop:k02:Bh}}{the proposition}}\label{app:a04:assembly}

Let $(\gamma,\beta,\sigma)$ be the exponents of \cref{prop:k01:Kpos}, $(\gamma',\beta',\sigma')$ those of \cref{hyp:k02:Kh}, and $(\gamma_b,\beta_b)$ those of \cref{app:a04:lem:BULK}. The diagonal pairs are taken with
\begin{equation}\label{app:a04:eq:alpha}
0<\alpha<\min\Big\{\tfrac12,\ \frac\gamma{\gamma+\beta},\ \frac{\gamma'}{\gamma'+\beta'},\ \frac{\gamma_b}{\gamma_b+\beta_b}\Big\};
\end{equation}
the Cauchy arguments of \cref{thm:k01:CA,thm:k02:CAh} hold for every $\alpha\in(0,\frac12)$ with $\alpha<\gamma/(\gamma+\beta)$, so the two further restrictions cost nothing, and one value of $\alpha$ is fixed once for the whole existence part. Let $\Sigma=\{(u_x,h_y):(x,y)\in Y\}$ be the compact family, as in \cref{app:a04:rows}.

\emph{Step 1 (the split).} Fix $k\ge k_0$, $j\defeq j(k)$ and $\sigma=(u_x,h_y)\in\Sigma$, and write $h\defeq h_y$. Both members of the diagonal pair live on the same unit lattice with the same indicator $\chi_j$, and \cref{app:a04:lem:split} applies to each. Hence, on the support of $\chi_j$,
\begin{align*}
S^{(j)}_{N,k}(\cdot;\sigma)+c_k(h)-S^{(j+1)}_{N,k+1}(\cdot;\sigma)-c_{k+1}(h)
&=\big[S^{(j)}_{N,k}(\cdot;(u_x,0))-S^{(j+1)}_{N,k+1}(\cdot;(u_x,0))\big]+\sum_{T\in\mathcal N_h}D[T]\\
&\quad+\Big(\big[\cV^{(j)}_k(h)+c_k(h)\big]-\big[\cV^{(j+1)}_{k+1}(h)+c_{k+1}(h)\big]\Big).
\end{align*}
At $\sigma=0$ the last two terms are absent and the left side is the left side of \eqref{eq:k01:Kpos} at $s=0$, $j=j(k)$.

\emph{Step 2 (rows K0 and SRC).} By \cref{app:a04:cor:rows} the first two terms are bounded by $\theta^{\mathrm{K0}}_k+\theta^{\mathrm{SRC}}_k$, uniformly on $\Sigma\cup\{0\}$.

\emph{Step 3 (row LOW).} The third term is field-independent and linear in $h$, and it vanishes at $h=0$. By \cref{app:a04:lem:loc} and \cref{app:a04:lem:alias} with $r=2$,
\[
\cV^{(j)}_k(h)+c_k(h)=\bar\beta^{(j)}_k\int h+E_k(h),\qquad
|E_k(h)|\le C\norm h_{L^1}\eps_k^{-\nu}e^{-c/\eta_k}+C\norm{\beta^{(j)}_k}_\infty\norm{\Delta^2h}_{L^1}\big(\eta_k^{4}+\eps_k^{4}\big),
\]
where $c_k(h)=z_k\sS_k(h)$ was used in the form of \eqref{app:a04:eq:loc}, and $|\bar\beta^{(j)}_k|\le\norm{\beta^{(j)}_k}_\infty$. The same holds for the $(k+1)$-model after $j+1$ steps, with $\eps_{k+1}=\eps_k/5$ and the same $\eta_k$. By \eqref{app:a04:eq:BULKinc} and \eqref{app:a04:eq:BULKsup}, the third term is bounded by
\[
\theta^{\mathrm{LOW}}_k(h)\defeq C\,5^{-\gamma_bj}\eta_k^{-\beta_b}\Big|\int h\Big|+C\norm h_{L^1}5^{\nu k}e^{-c5^{\alpha k-1}}+C\norm{\Delta^2h}_{L^1}\big(\eta_k^{2}+\eta_k^{-2}\eps_k^{4}\big).
\]
Its coefficients are linear in $(|\int h|,\norm h_{L^1},\norm{\Delta^2h}_{L^1})$, which are bounded on $\{h_y:y\in Y\}$; so $\theta^{\mathrm{LOW}}_k(h_y)\le\theta^{\mathrm{LOW}}_k$ with $\theta^{\mathrm{LOW}}_k$ independent of $\sigma\in\Sigma$, and $\theta^{\mathrm{LOW}}_k(0)=0$.

\emph{Step 4 (geometric decay).} Along the diagonal pairs $\eta_k\le5^{-\alpha k}$, $j(k)\ge(1-\alpha)k-1$ and $k-j(k)\le\alpha k+1$ (Step~6 of the proof of \cref{thm:k01:CA}). Hence $\eta_k^{\sigma}$, $\eta_k^{\sigma'}$, $\eta_k^{2}$ and $\eta_k^{-2}\eps_k^{4}\le5^{2(\alpha k+1)-4k}$ decay geometrically; $5^{\nu k}e^{-c5^{\alpha k-1}}$ decays faster than any geometric sequence; and for $(\gamma_*,\beta_*)\in\{(\gamma,\beta),(\gamma',\beta'),(\gamma_b,\beta_b)\}$,
\[
5^{-\gamma_*j(k)}\eta_k^{-\beta_*}\le5^{\gamma_*+\beta_*}\,5^{-(\gamma_*(1-\alpha)-\beta_*\alpha)k},
\]
which decays geometrically because $\alpha<\gamma_*/(\gamma_*+\beta_*)$ by \eqref{app:a04:eq:alpha}. Therefore
\[
\theta_k\defeq\theta^{\mathrm{K0}}_k+\theta^{\mathrm{SRC}}_k+\theta^{\mathrm{LOW}}_k
\]
decays geometrically, does not depend on $\sigma\in\Sigma$, and by Steps 1--3 the left side of \eqref{eq:k02:B} is bounded on the support of $\chi_{j(k)}$ by $\theta_k(\sigma)\defeq\theta^{\mathrm{K0}}_k+\theta^{\mathrm{SRC}}_k+\theta^{\mathrm{LOW}}_k(h_\sigma)\le\theta_k$ for $\sigma\in\Sigma$, and by $\theta_k(0)\defeq\theta^{\mathrm{K0}}_k$ at $\sigma=0$. This is \cref{prop:k02:Bh}. \qed

\begin{remark}
The proof uses no convergence and no bound of the vacuum constants $v_{k,j}$, no bound on $\Zk(s)$, and no property of the scheme $z$ beyond $|z_k|\le C(1+k)$ (\cref{lem:k02:cnum}(a)) and the identity $\beta^{(j)}_k=v^{(j)}_{k,\infty}+z_k$ of \cref{app:a04:def:bulk}: the $c$-number $c_k(h)=z_k\sS_k(h)$ is the one for which $\cV^{(j)}_k(h)+c_k(h)$ is, up to the errors of \cref{app:a04:lem:loc,app:a04:lem:alias}, the block mean of the bulk density times $\int h$, and the bulk density is the difference between the low-order one-point function of the renormalised Wick square in the full lattice model and in King's fluctuation measure.
\end{remark}
